\label{chap:log-ht-monodromy}%
\label{conv:five-levels-log-ht}%


\providecommand{\A}{\mathcal A}
\providecommand{\fS}{\mathfrak{S}}
\providecommand{\B}{\mathcal B}
\providecommand{\C}{\mathcal C}
\providecommand{\D}{\mathcal D}
\providecommand{\E}{\mathcal E}
\providecommand{\F}{\mathcal F}
\providecommand{\FM}{\mathrm{FM}}
\providecommand{\Conf}{\mathrm{Conf}}
\providecommand{\HT}{\mathrm{HT}}
\providecommand{\MC}{\mathrm{MC}}
\providecommand{\End}{\mathrm{End}}
\providecommand{\Hom}{\mathrm{Hom}}
\providecommand{\id}{\mathrm{id}}
\providecommand{\im}{\mathrm{im}}
\providecommand{\Res}{\mathrm{Res}}
\providecommand{\Aut}{\mathrm{Aut}}
\providecommand{\PP}{\mathbb P}
\providecommand{\AA}{\mathbb A}
\providecommand{\CC}{\mathbb C}
\providecommand{\RR}{\mathbb R}
\providecommand{\ZZ}{\mathbb Z}
\providecommand{\kk}{\Bbbk}
\providecommand{\h}{\hbar}
\providecommand{\ot}{\otimes}
\providecommand{\wot}{\widehat\otimes}
\providecommand{\s}{\mathrm s}
\providecommand{\ds}{d_{\mathrm{dR}}}
\providecommand{\CYB}{\mathrm{CYB}}
\providecommand{\PT}{\mathrm{PT}}
\providecommand{\Disk}{\mathrm{Disk}}
\providecommand{\SN}{\boldsymbol{\nabla}}
\providecommand{\bboxprod}{\mathop{\boxtimes}\displaylimits}

\section{Introduction: the local law, the global law, and the bridge}

The local algebra of a chiral or holomorphic-topological theory is not its global geometry. Local operators know how fields collide; the global theory knows how states transport when collisions are avoided. Between the two sits configuration-space geometry: blown up until collisions become divisors, regularized until graph integrals become logarithmic forms, reduced until only degree-zero states remain. This geometry converts local algebra into monodromy.

The guiding principle of this manuscript is therefore simple:
\begin{quote}
\emph{the local law is algebraic, the global law is monodromic, and the bridge is logarithmic collision geometry.}
\end{quote}

From this viewpoint, the subjects in play (chiral and factorization algebra, $A_\infty$ deformation theory, dg-shifted Yangians, Fulton-MacPherson compactifications, perturbative holomorphic-topological field theory, logarithmic geometry, pure braid groups, periods, and spectral/chromatic stabilization) become successive aspects of one mechanism.
\begin{enumerate}[label=\textup{(\arabic*)}]
 \item A binary collision kernel $r(z)$ satisfying a Yang-Baxter identity defines a logarithmic flat connection on configuration space.
 \item Replacing strict multiplication by an $A_\infty$ structure forces the connection to lift to a bar-valued superconnection; its curvature is the Maurer-Cartan defect.
 \item In interacting holomorphic-topological theories, the superconnection is realized analytically by renormalized compactified graph integrals on mixed holomorphic/topological configuration spaces.
 \item When the internal line complex is classically resolved in degree $0$, the superconnection reduces canonically to an ordinary logarithmic connection on reduced states.
 \item Regularized transport between tangential collision zones then produces an associator, a braiding, and hence a monodromic braided tensor structure.
\end{enumerate}

The mechanism is packaged into a single synthesis theorem. Whenever a result depends on analytic or perturbative inputs external to the manuscript, we say so explicitly.

\begin{theorem}[Synthesis theorem; \ClaimStatusProvedHere]\label{thm:synthesis}
Let $\kk$ be a field of characteristic zero.
\begin{enumerate}[label=\textup{(\roman*)}]
 \item A strict rational dg-shifted Yangian $(Y,r(z),T)$ determines a logarithmic flat shifted KZ/FM connection on $\Conf_n(\AA^1)$ whose residues extend to $\FM_n(\PP^1\mid\infty)$ and satisfy infinitesimal braid relations.
 \item Support-indexed collision fields $\{D_S\}_{\varnothing\neq S\subset[n]}$ in an $A_\infty$ algebra determine a bar-valued superconnection
 \[
 \SN_n=\ds+\widehat\Theta_n+b_n+I_{\Gamma_n},\qquad \Gamma_n=\sum_{\varnothing\neq S\subset[n]}D_S,
 \]
 whose curvature is
 \[
 \SN_n^2 = I_{(\ds+\Theta_n)\Gamma_n+\MC(\Gamma_n)}.
 \]
 Hence flatness is equivalent to a supportwise Maurer-Cartan hierarchy and yields a recursive obstruction theory for higher collision homotopies.
 \item In a quasi-linear interacting three-dimensional holomorphic-topological theory with anomaly-free perturbation theory, renormalized compactified graph integrals on mixed logarithmic configuration spaces produce an analytic logarithmic superconnection of the same form, and its flatness is a compactified-Stokes identity.
 \item If the tree-level line complex is resolved in degree $0$, then the full quantum line complex remains concentrated in degree $0$, admits a canonical $\h$-adic strong deformation retract, and the reduced superconnection collapses to an ordinary logarithmic flat connection
 \[
 \A_n = p_n A_n^{(1)} i_n.
 \]
 In particular, the reduced monodromic theory is controlled directly by the projected renormalized logarithmic $1$-form sector.
 \item Regularized transport of the reduced connection between tangential collision zones produces an associator, a braiding, pentagon and hexagon identities, and pure braid group representations on reduced line states.
 \item After spectralization, smashing localizations commute with factorization homology. Consequently periodicity and chromatic localization are compatible with the local-to-global passage.
\end{enumerate}
\end{theorem}

Every clause of Theorem~\ref{thm:synthesis} is proved later in the article, except where we explicitly mark conjectural bridges beyond the present theorem package.

\subsection*{Guiding references and scope}
The formal development is self-contained, but builds on: Beilinson--Drinfeld's chiral/factorization viewpoint \cite{BD04}; Ayala--Francis on factorization homology \cite{AF15}; Costello--Dimofte--Gaiotto on holomorphic twists and boundary chiral algebras \cite{CDG20}; Gaiotto--Kulp--Wu on higher operations and compactified configuration-space regularization \cite{GKW25}; Khan--Zeng on three-dimensional Poisson-vertex gauge theory \cite{KZ25}; Mok on logarithmic Fulton-MacPherson spaces \cite{Mok25}; Dimofte--Niu--Py on line operators and dg-shifted Yangians \cite{DNP25}; and Latyntsev on factorisation quantum groups \cite{Lat23}.

\begin{remark}[Log-geometric versus analytic Fulton--MacPherson; scope]\label{rem:log-geometric-vs-analytic-FM}
Two distinct compactifications of $\Conf_n$ appear throughout this chapter, and they should not be conflated. As a scope guard: \emph{logarithmic HT monodromy} refers to log-FM geometry, \emph{not} logarithmic VOAs. Within the log-FM framework itself we record the following algebraic/analytic split:
\begin{enumerate}[label=\textup{(\alph*)}]
 \item \textbf{log-FM} $=$ $\FM_n^{\mathrm{KN}}$, the Kato--Nakayama log-scheme Fulton--MacPherson compactification carrying the fs log-structure along the collision boundary divisors. This incarnation is \emph{algebraic/motivic}: sheaves of log-differentials $\Omega^\bullet_{\FM_n^{\mathrm{KN}}}(\log)$, Gauss--Manin connections on log de Rham cohomology, and residue morphisms on boundary strata are all available.
 \item \textbf{analytic-FM} $=$ $\FM_n^{\mathrm{an}}$, the complex-analytic compactification obtained by polar blow-up of collision loci; its boundary carries \emph{Stokes-regularity data} (tangential base points, real oriented blow-up, sectorial decompositions) needed for the classification of irregular-singular flat connections.
 \item \textbf{Bridge.} The Kato--Nakayama comparison $\FM_n^{\mathrm{KN}} \simeq \FM_n^{\mathrm{an}}$ is a quasi-equivalence in the $\infty$-categorical sense (agreement on log de Rham / Betti cohomology after Kato--Nakayama realization); at the cohomological level the two frameworks compute the same periods, but at the chain/datum level log-FM carries algebraic structure while analytic-FM carries wild-monodromy Stokes data.
\end{enumerate}
\emph{Programme usage.} The irregular-singular KZB composition theorem (\cref{thm:irregular-kzb-composition-generic-level} in the modular operad package) is formulated on $\FM_n^{\mathrm{an}}$: its content is a Stokes classification at boundary divisors where the KZB connection is irregular at generic non-integral level. The modular bootstrap $H^2=0$ vanishing (\cref{thm:mb-H2-vanishing}) is formulated on $\FM_n^{\mathrm{KN}}$: its content is Gauss--Manin flatness of the log de Rham cohomology of the log-FM stratification. The two statements meet at cohomological level via the Kato--Nakayama comparison; one does not substitute for the other at the chain level.
\end{remark}

\section{First principles and conventions}\label{sec:log-ht-foundations}

We work over a characteristic-zero ground field $\kk$. Unless stated otherwise, grading conventions are cohomological and tensor products are completed when necessary.

\subsection{Suspended bar conventions}
Let $A$ be an $A_\infty$ algebra. We use the suspended convention: the structure maps are
\[
 m_k : (\s A)^{\otimes k} \longrightarrow \s A
\]
of cohomological degree $+1$. Write $\bar{A}:=\ker(\epsilon)$ for the augmentation ideal. The completed tensor coalgebra
\[
 \mathrm B(A):=T^c(\s \bar{A})=\prod_{m\ge 0}(\s \bar{A})^{\otimes m}
\]
comes equipped with a coderivation $b$ of degree $+1$ satisfying $b^2=0$.

For $x\in A^1$, define the insertion coderivation $I_x$ on $\mathrm B(A)$ by
\[
 I_x(\s a_1\otimes\cdots\otimes \s a_m)
 :=
 \sum_{j=0}^{m}
 \s a_1\otimes\cdots\otimes\s a_j\otimes\s x\otimes\s a_{j+1}\otimes\cdots\otimes\s a_m.
\]
The Maurer-Cartan polynomial is
\[
 \MC(x):=\sum_{\ell\ge 1} m_\ell(x,\dots,x).
\]

\begin{lemma}[Bar insertion identity; \ClaimStatusProvedHere]\label{lem:bar-insertion}
For every degree-one element $x\in A^1$ one has
\[
 [b,I_x]+I_x^2 = I_{\MC(x)}.
\]
\end{lemma}

\begin{proof}
Both sides are coderivations, so it suffices to compare corestrictions to the cogenerators $\s \bar{A}$. The corestriction of $[b,I_x]$ records all terms in which a single higher multiplication $m_\ell$ acts on a bar word containing exactly one inserted copy of $x$; the corestriction of $I_x^2$ records the contributions in which two inserted copies of $x$ become adjacent and are then absorbed into the same higher multiplication. Collecting all such possibilities yields precisely the full Maurer-Cartan polynomial $\MC(x)$. This is the suspended bar-coalgebra form of the usual $A_\infty$ Maurer-Cartan identity.
\end{proof}

\subsection{Configuration spaces and logarithmic Fulton-MacPherson geometry}
For a smooth curve $C$ and a divisor $D\subset C$, write
\[
 \Conf_n(C\setminus D)=\{(z_1,\dots,z_n)\in (C\setminus D)^n\mid z_i\neq z_j\text{ for }i\neq j\}
\]
for the ordered configuration space. We set
\[
 X_n:=\FM_n(\PP^1\mid\infty)
\]
for Mok's logarithmic Fulton-MacPherson compactification of the pair $(\PP^1,\infty)$ \cite{Mok25}. Near a boundary stratum indexed by a collision tree $T$, there exist local scale coordinates $t_e$ indexed by internal edges $e\in E(T)$ such that pairwise differences factor as
\[
 z_i-z_j = \Big(\prod_{e\in P(i,j)} t_e\Big)u_{ij},
\]
where $u_{ij}$ is a unit and $P(i,j)$ is the set of tree edges separating $i$ from $j$. Thus collisions become normal-crossings divisors in logarithmic coordinates.

\subsection{Strict rational dg-shifted Yangians}
The full notion of dg-shifted Yangian belongs to a broader program \cite{DNP25}. For the strict rational theory developed below we isolate only the structure used.

\begin{definition}[Strict rational dg-shifted Yangian]\label{def:strict-yangian}
A \emph{strict rational dg-shifted Yangian} consists of the following data:
\begin{enumerate}[label=\textup{(\alph*)}]
 \item an associative dg algebra $Y$;
 \item a degree-zero derivation $T$ (the translation operator);
 \item a meromorphic element $r(z)\in (Y\otimes Y)(z)$, regular at infinity, with a simple pole at the origin,
 \[
 r(z)=\frac{\Omega}{z}+r_{\mathrm{reg}}(z),\qquad \Omega\in Y\otimes Y;
 \]
 The residue element absorbs
 the level $k$, so the affine realization reads
 $r(z)=k\,\Omega_{\mathrm{aff}}/z$ with $\Omega_{\mathrm{aff}}$ the
 quadratic Casimir tensor of $\mathfrak g$, vanishing at $k=0$, cf.\
 the affine specialization at the start of Section~\ref{sec:strict}.)
 \item the parameter-dependent classical Yang-Baxter equation
 \[
 [r_{12}(u),r_{13}(u+v)] + [r_{12}(u),r_{23}(v)] + [r_{13}(u+v),r_{23}(v)] = 0.
 \]
\end{enumerate}
\end{definition}

The datum $(Y,r,T)$ is the strict shadow of the full $A_\infty$ structure treated later.

\section{The strict logarithmic theory: the shifted KZ/FM connection}\label{sec:strict}

We begin with the strict theory because it exposes the geometric heart of the subject with maximal clarity: pairwise collisions, logarithmic poles, and Yang-Baxter flatness.

\subsection{Construction}
Let $(Y,r,T)$ be a strict rational dg-shifted Yangian and let $V_1,\dots,V_n$ be strict left $Y$-modules. Write $z_{ij}:=z_i-z_j$. On the trivial vector bundle over $\Conf_n(\AA^1)$ with fiber $V_1\otimes\cdots\otimes V_n$, define
\[
 \nabla_n^Y
 :=
 d-\sum_{1\le i<j\le n}\rho_{ij}(r(z_{ij}))\,d(z_{ij}),
\]
where $\rho_{ij}$ acts in factors $i$ and $j$ and trivially elsewhere.

\begin{definition}[Shifted KZ/FM connection]
The connection $\nabla_n^Y$ is called the \emph{shifted KZ/FM connection} attached to $Y$.
\end{definition}

When $r(z)=k\,\Omega/z$ (the affine collision residue at level~$k$), this is exactly the classical rational KZ connection.

\subsection{Flatness and residues}

\begin{theorem}[Flatness of the shifted KZ/FM connection; \ClaimStatusProvedHere]\label{thm:strict-flatness}
For every strict rational dg-shifted Yangian, the shifted KZ/FM connection is flat:
\[
 (\nabla_n^Y)^2=0.
\]
\end{theorem}

\begin{proof}
Set
\[
 A:=\sum_{i<j}\rho_{ij}(r(z_{ij}))\,d(z_{ij}).
\]
Then $\nabla_n^Y=d-A$, so its curvature is $F=dA-A\wedge A$. First, $dA=0$, because every summand has the form $f(z_{ij})d(z_{ij})$ and hence
\[
 d\big(f(z_{ij})d(z_{ij})\big)=f'(z_{ij})d(z_{ij})\wedge d(z_{ij})=0.
\]
Thus only $A\wedge A$ remains.

Terms supported on disjoint pairs commute, so only triples $i<j<k$ contribute. For one such triple,
\begin{align*}
 A\wedge A\big|_{ijk}
 &= [r_{ij}(z_{ij}),r_{ik}(z_{ik})]dz_{ij}\wedge dz_{ik}
 + [r_{ij}(z_{ij}),r_{jk}(z_{jk})]dz_{ij}\wedge dz_{jk} \\
 &\qquad + [r_{ik}(z_{ik}),r_{jk}(z_{jk})]dz_{ik}\wedge dz_{jk}.
\end{align*}
Since $z_{ik}=z_{ij}+z_{jk}$, we also have
\[
 dz_{ij}\wedge dz_{ik}=dz_{ij}\wedge dz_{jk},
 \qquad
 dz_{ik}\wedge dz_{jk}=dz_{ij}\wedge dz_{jk}.
\]
Hence the triple contribution equals
\[
 \Big([r_{12}(u),r_{13}(u+v)] + [r_{12}(u),r_{23}(v)] + [r_{13}(u+v),r_{23}(v)]\Big)_{u=z_{ij},\,v=z_{jk}}dz_{ij}\wedge dz_{jk},
\]
which vanishes by the classical Yang-Baxter equation. Summing over all triples gives $A\wedge A=0$, hence $F=0$.
\end{proof}

\begin{lemma}[Residues and infinitesimal braid operators; \ClaimStatusProvedHere]\label{lem:inf-braid}
If $r(z)=k\,\Omega/z+O(1)$ at level~$k$ near $z=0$ and $\Omega_{ij}$ denotes the action of $\Omega$ in factors $i,j$, then
\[
 [\Omega_{ij},\Omega_{kl}]=0 \qquad \text{if }\{i,j\}\cap\{k,l\}=\varnothing,
\]
and
\[
 [\Omega_{ij},\Omega_{ik}+\Omega_{jk}]=0
\]
for distinct $i,j,k$.
\end{lemma}

\begin{proof}
The disjoint-support commutators are tautologically zero. For the three-term relation, substitute $u=\lambda x$ and $v=\lambda y$ into the Yang-Baxter equation, multiply by $\lambda^2$, and let $\lambda\to 0$. Using
\[
 r(\lambda x)=\frac{k\,\Omega}{\lambda x}+O(1),
 \qquad
 r(\lambda(x+y))=\frac{k\,\Omega}{\lambda(x+y)}+O(1),
\]
one obtains
\[
 \frac{[\Omega_{12},\Omega_{13}]}{x(x+y)}+
 \frac{[\Omega_{12},\Omega_{23}]}{xy}+
 \frac{[\Omega_{13},\Omega_{23}]}{(x+y)y}=0.
\]
Multiplying by $xy(x+y)$ and comparing the coefficients of $x$ and $y$ yields the claim.
\end{proof}

\begin{proposition}[Logarithmic extension and edge residues; \ClaimStatusProvedHere]\label{prop:strict-log-ext}
Assume $r(z)$ is regular at infinity. Then the shifted KZ/FM connection extends to a logarithmic connection on
\[
 X_n=\FM_n(\PP^1\mid\infty).
\]
Near a boundary stratum indexed by a collision tree $T$, there are local scale coordinates $t_e$ such that
\[
 \nabla_n^Y = d-\sum_{e\in E(T)} \Omega_e\,\frac{dt_e}{t_e}+\text{regular terms},
\]
where
\[
 \Omega_e:=\sum_{i<j:\,e\in P(i,j)}\Omega_{ij}.
\]
\end{proposition}

\begin{proof}
In a tree chart one has
\[
 z_i-z_j = \Big(\prod_{e\in P(i,j)}t_e\Big)u_{ij},
\]
with $u_{ij}$ a unit. Since $r(z)=k\,\Omega/z+O(1)$ at level~$k$ one has
\[
 r(z_i-z_j)d(z_i-z_j)=k\,\Omega\,d\log(z_i-z_j)+\text{regular terms}.
\]
Expanding the logarithmic derivative gives
\[
 d\log(z_i-z_j)=\sum_{e\in P(i,j)}\frac{dt_e}{t_e}+d\log(u_{ij}).
\]
Summing over all pairs $i<j$ therefore yields
\[
 \sum_{i<j}r(z_i-z_j)d(z_i-z_j)=k\sum_{e\in E(T)}\Big(\sum_{i<j:\,e\in P(i,j)}\Omega_{ij}\Big)\frac{dt_e}{t_e}+\text{regular terms},
\]
which is the claimed logarithmic form \textup{(}with overall level factor~$k$; at $k=0$ the connection reduces to the exterior derivative, as required\textup{)}.
\end{proof}

\begin{corollary}[Local monodromy around a boundary divisor; \ClaimStatusProvedHere]\label{cor:strict-monodromy}
The local monodromy around the divisor $t_e=0$ is conjugate to
\[
 \exp(2\pi i\,\Omega_e).
\]
\end{corollary}

\begin{proof}
After a local holomorphic gauge transformation removing the regular part, the connection has the model form $d-\Omega_e\,dt_e/t_e$ in the normal direction to the divisor. Parallel transport around $t_e=\varepsilon e^{2\pi i s}$ gives the stated exponential.
\end{proof}

\begin{remark}[Enumerative and arithmetic shadows of the strict theory]
The boundary of $X_n$ is indexed by collision trees. Proposition~\ref{prop:strict-log-ext} decorates every internal edge of every collision tree with an explicit residue operator $\Omega_e$. Thus the compactification is not only combinatorial; it is a tree skeleton weighted by operator-valued residues. If the algebra and modules are defined over a subfield $K\subset\CC$, the connection and all its residues are defined over $K$, producing a logarithmic local system with arithmetic structure.
\end{remark}

\section{The full $A_\infty$ collision theory}\label{sec:ainfty}

The strict connection sees only binary collisions. The genuinely derived theory must also remember higher collisions, and these are governed by higher homotopies. The correct home for them is not an ordinary vector bundle but the completed bar coalgebra.

\subsection{Support-indexed collision fields}
Let $A$ be an $A_\infty$ algebra in the suspended convention, endowed with a degree-zero translation derivation $T$. Let $V_1,\dots,V_n$ be line-state objects and set
\[
 E_n:=\End(V_1\otimes\cdots\otimes V_n)\wot A.
\]
For every nonempty subset $S\subset[n]$, let
\[
 D_S(\mathbf z_S)\in E_n^1
\]
be a degree-one element supported exactly on the tensor factors indexed by $S$. The intended interpretation is
\[
 D_{\{i\}}=\mu_i(z_i),
 \qquad
 D_{\{i,j\}}=r_{ij}(z_i-z_j),
 \qquad
 D_S=h_S(\mathbf z_S)\quad (|S|\ge 3),
\]
where $\mu_i$ is the translated one-line Maurer-Cartan coupling, $r_{ij}$ is the binary collision kernel, and $h_S$ is a genuinely higher collision correction.

Define the total collision field
\[
 \Gamma_n(\mathbf z):=\sum_{\varnothing\neq S\subset[n]}D_S(\mathbf z_S).
\]
Let $\mathrm B(E_n)$ be the completed bar coalgebra of $E_n$, with bar coderivation $b_n$ induced by the $A_\infty$ structure. Let
\[
 \Theta_n:=\sum_{i=1}^n dz_i\,T^{(i)}
\]
be the translation one-form acting on $E_n$, and let $\widehat\Theta_n$ be the induced coderivation on $\mathrm B(E_n)$.

\begin{definition}[Full $A_\infty$ superconnection]\label{def:superconnection}
The \emph{full $A_\infty$ collision superconnection} is
\[
 \SN_n := \ds + \widehat\Theta_n + b_n + I_{\Gamma_n}.
\]
\end{definition}

\subsection{The master curvature formula}

\begin{theorem}[Master curvature formula; \ClaimStatusProvedHere]\label{thm:master-curvature}
The curvature of the full $A_\infty$ collision superconnection is again an insertion coderivation:
\[
 \SN_n^2 = I_{\mathcal K_n},
\]
where
\[
 \mathcal K_n := (\ds+\Theta_n)\Gamma_n + \MC(\Gamma_n)
 = (\ds+\Theta_n)\Gamma_n + \sum_{\ell\ge 1} m_\ell(\Gamma_n,\dots,\Gamma_n).
\]
In particular,
\[
 \SN_n^2=0
 \quad\Longleftrightarrow\quad
 (\ds+\Theta_n)\Gamma_n+\MC(\Gamma_n)=0.
\]
\end{theorem}

\begin{proof}
Since $\ds^2=0$, $\widehat\Theta_n^2=0$, and $b_n^2=0$, only mixed commutators with $I_{\Gamma_n}$ remain. First,
\[
 [\ds+\widehat\Theta_n,I_{\Gamma_n}] = I_{(\ds+\Theta_n)\Gamma_n},
\]
because de Rham differentiation and translation act coefficientwise on the inserted element.

Second, by Lemma~\ref{lem:bar-insertion},
\[
 [b_n,I_{\Gamma_n}] + I_{\Gamma_n}^2 = I_{\MC(\Gamma_n)}.
\]
Adding the identities yields
\[
 \SN_n^2 = I_{(\ds+\Theta_n)\Gamma_n+\MC(\Gamma_n)},
\]
which is the claim.
\end{proof}

The curvature is the Maurer--Cartan defect of the total collision field after transport by de Rham and translation.

\subsection{Supportwise curvature hierarchy}
For every nonempty subset $S\subset[n]$, let $(\mathcal K_n)_S$ denote the support-$S$ component of $\mathcal K_n$. Since the support of
\[
 m_\ell(D_{S_1},\dots,D_{S_\ell})
\]
is $S_1\cup\cdots\cup S_\ell$, the flatness equation decomposes supportwise.

\begin{theorem}[Support hierarchy; \ClaimStatusProvedHere]\label{thm:support-hierarchy}
The flatness equation $\SN_n^2=0$ is equivalent to the family of equations
\[
 (\ds+\Theta_n)D_S +
 \sum_{\ell\ge 1}
 \sum_{\substack{S_1,\dots,S_\ell\neq\varnothing \\
 S_1\cup\cdots\cup S_\ell = S}}
 m_\ell(D_{S_1},\dots,D_{S_\ell})
 =0
\]
for all nonempty $S\subset[n]$.
\end{theorem}

\begin{proof}
Take the support-$S$ component of the identity
\[
 (\ds+\Theta_n)\Gamma_n+\MC(\Gamma_n)=0.
\]
The support bookkeeping is exactly as stated.
\end{proof}

There is an equivalent Möbius-inversion form. For $R\subset[n]$ set
\[
 \Gamma_R:=\sum_{\varnothing\neq U\subseteq R}D_U.
\]
Then the support-$S$ equation is equivalent to
\[
 (\ds+\Theta_n)D_S + \sum_{\varnothing\neq R\subseteq S}(-1)^{|S|-|R|}\MC(\Gamma_R)=0.
\]
The support hierarchy therefore behaves like an inclusion-exclusion principle for collision clusters.

\subsection{The first three levels}
The hierarchy is already highly informative in low support.

\begin{corollary}[The first three curvature layers; \ClaimStatusProvedHere]\label{cor:first-three}
The support hierarchy specializes to the following equations.
\begin{enumerate}[label=\textup{(\roman*)}]
 \item For $S=\{i\}$,
 \[
 (\ds+\Theta_n)\mu_i + \MC(\mu_i)=0.
 \]
 \item For $S=\{i,j\}$,
 \[
 (\ds+\Theta_n)r_{ij} + \MC(\mu_i+\mu_j+r_{ij}) - \MC(\mu_i)-\MC(\mu_j)=0.
 \]
 \item For $S=\{i,j,k\}$, writing
 \[
 \Gamma_{ij}:=\mu_i+\mu_j+r_{ij},
 \qquad
 \Gamma_{ijk}:=\mu_i+\mu_j+\mu_k+r_{ij}+r_{ik}+r_{jk}+h_{ijk},
 \]
 one has
 \[
 (\ds+\Theta_n)h_{ijk} + \MC(\Gamma_{ijk})-\MC(\Gamma_{ij})-\MC(\Gamma_{ik})-\MC(\Gamma_{jk})+\MC(\mu_i)+\MC(\mu_j)+\MC(\mu_k)=0.
 \]
\end{enumerate}
\end{corollary}

If the singleton equations hold, the triple equation simplifies to
\[
 (\ds+\Theta_n)h_{ijk} + \MC(\Gamma_{ijk})-\MC(\Gamma_{ij})-\MC(\Gamma_{ik})-\MC(\Gamma_{jk})=0.
\]
Thus the classical $A_\infty$ Yang-Baxter defect is not the end of the story; it is the source term for a triple collision homotopy.

\subsection{Obstruction theory}

\begin{definition}[Twisted differential at the binary stage]
Let
\[
 \Gamma_{\le 2}:=\sum_{|S|\le 2}D_S.
\]
The \emph{binary-stage twisted differential} on support-$\ge 3$ data is
\[
 \delta_{\le 2}(X)
 :=
 (\ds+\Theta_n)X + \sum_{p,q\ge 0} m_{p+1+q}(\Gamma_{\le 2}^{\otimes p},X,\Gamma_{\le 2}^{\otimes q}).
\]
\end{definition}

\begin{theorem}[First obstruction theorem; \ClaimStatusProvedHere]\label{thm:first-obstruction}
Assume that the support-$1$ and support-$2$ equations hold. Let
\[
 \Omega_{ijk}:=\Big((\ds+\Theta_n)\Gamma_{\le 2}+\MC(\Gamma_{\le 2})\Big)_{\{i,j,k\}}.
\]
Then
\[
 \delta_{\le 2}(\Omega_{ijk})=0.
\]
Hence $[\Omega_{ijk}]\in H^2(\delta_{\le 2})$ is the obstruction class to the existence of a triple correction $h_{ijk}$ solving the support-$3$ flatness equation.
\end{theorem}

\begin{proof}
Let
\[
 \SN_{\le 2}:=\ds+\widehat\Theta_n+b_n+I_{\Gamma_{\le 2}}.
\]
By Theorem~\ref{thm:master-curvature},
\[
 \SN_{\le 2}^2 = I_{\Omega},
\]
where $\Omega$ has support at least $3$ and its support-$3$ component is precisely $\Omega_{ijk}$. Apply the Bianchi identity:
\[
 [\SN_{\le 2},\SN_{\le 2}^2]=0.
\]
Taking the support-$\{i,j,k\}$ corestriction to cogenerators gives exactly
\[
 \delta_{\le 2}(\Omega_{ijk})=0.
\]
The support-$3$ flatness equation is the inhomogeneous equation
\[
 \delta_{\le 2}(h_{ijk})=-\Omega_{ijk}
\]
up to terms of higher order in any auxiliary pronilpotent filtration. Thus $h_{ijk}$ exists if and only if the class of $\Omega_{ijk}$ vanishes.
\end{proof}

\begin{corollary}[Recursive support filtration; \ClaimStatusProvedHere]\label{cor:recursive-support}
Fix an auxiliary pronilpotent filtration in which nonlinear terms are of strictly higher order. Suppose support-$<m$ data have been constructed. Then for each $|S|=m$, the defect
\[
 \mathfrak O_S:=\Big((\ds+\Theta_n)\Gamma_{<m}+\MC(\Gamma_{<m})\Big)_S
\]
is closed for the twisted differential determined by $\Gamma_{<m}$. Its cohomology class is the obstruction to adjoining a support-$S$ correction $D_S$. If all these obstruction groups vanish, binary data extend recursively to a full flat collision superconnection.
\end{corollary}

\begin{remark}[The formal moral]
The $A_\infty$ Yang-Baxter relation is not a single equation but the first nontrivial face of an infinite supportwise curvature tower. Binary collisions generate curvature in support $3$; support-$3$ homotopies absorb that curvature and move the defect to support $4$; and so on. The full theory is therefore intrinsically recursive and operadic.
\end{remark}

\section{Analytic realization in interacting holomorphic-topological theory}\label{sec:analytic}

The formal theory of the previous section becomes analytic in a concrete interacting three-dimensional holomorphic-topological field theory. Compactified Stokes theory on the relevant configuration spaces realizes the formal Maurer--Cartan equations as literal boundary identities of renormalized graph integrals.

\subsection{The quasi-linear holomorphic-topological setting}
We work with a perturbative three-dimensional holomorphic-topological theory on
\[
 M=\CC_z\times \RR_t
\]
in the quasi-linear class highlighted in \cite{DNP25}. Concretely, this includes holomorphic-topological twists of many $3$d $\mathcal N=2$ gauge theories with arbitrary Chern-Simons level, linear matter, and superpotential, under the assumptions for which the perturbative line OPE is exactly resumable. Let $L_1,\dots,L_n$ be perturbative line operators supported on parallel copies of the topological direction, with state spaces $V_1,\dots,V_n$.

The analytic data we need are the graph-valued forms computing the higher operations of the theory in the sense of \cite{GKW25,DNP25}. The issue is that these forms live on singular configuration spaces. The cure is to compactify and logarithmically resolve the relevant collision strata.

\subsection{Mixed logarithmic configuration spaces}
Fix external holomorphic positions
\[
 \mathbf z=(z_1,\dots,z_n)\in \Conf_n(\CC),
 \qquad
 D_{\mathbf z}:=\{z_1,\dots,z_n,\infty\}\subset \PP^1.
\]
For a connected admissible graph $\Gamma$, let $V_b(\Gamma)$ denote the set of internal bulk vertices and $T(\Gamma)$ the set of all time variables carried by internal bulk and line vertices. The raw configuration space is
\[
 C_\Gamma(\mathbf z)
 :=
 \Conf_{V_b(\Gamma)}(\PP^1\setminus D_{\mathbf z})
 \times
 \Conf_{T(\Gamma)}^{\mathrm{ord}}(\RR),
\]
where the second factor is the chamber of ordered time variables compatible with line ordering.

\begin{definition}[Mixed logarithmic compactification]\label{def:mixed-log-config}
The \emph{mixed logarithmic configuration space} associated to $\Gamma$ and $\mathbf z$ is the iterated clean blowup
\[
 \mathfrak C_\Gamma(\mathbf z)
 :=
 \Big(
 \FM_{V_b(\Gamma)}(\PP^1\mid D_{\mathbf z})
 \times
 \overline{\Conf}^{\mathrm{ord}}_{T(\Gamma)}(\RR)
 \Big)^{\mathrm{mix}},
\]
obtained by further blowing up all mixed incidence strata corresponding to connected divergent subgraphs: bulk-bulk collisions, approach of bulk vertices to the divisor $D_{\mathbf z}$, time collisions, and their clean intersections forced by a connected collision pattern.
\end{definition}

The holomorphic and topological singularities are not independent: a divergent subgraph can force simultaneous holomorphic and temporal scale collapse, and those loci must be blown up together.

\begin{proposition}[Local coordinates and log smoothness; \ClaimStatusProvedHere]\label{prop:mixed-log-smooth}
The mixed logarithmic configuration space $\mathfrak C_\Gamma(\mathbf z)$ carries a natural fs logarithmic structure and is log smooth over $X_n=\FM_n(\PP^1\mid\infty)$. Its codimension-one boundary divisors are indexed by connected mixed collision types. Near a boundary stratum indexed by a mixed collision forest $F$, there are local scale coordinates
\[
 r_e\ge 0 \quad (e\in E_h(F)),
 \qquad
 s_f\ge 0 \quad (f\in E_t(F)),
\]
and unit coordinates $\xi,\tau$ such that
\[
 w_a-w_b = \Big(\prod_{e\in P_h(a,b)}r_e\Big)\xi_{ab},
 \qquad
 t_u-t_v = \Big(\prod_{f\in P_t(u,v)}s_f\Big)\tau_{uv}.
\]
\end{proposition}

\begin{proof}
Mok's logarithmic Fulton-MacPherson spaces are log smooth, and the ordered real compactification of ordered points on a line is a manifold with corners. The mixed incidence strata we blow up are clean intersections of boundary strata indexed by connected collision subgraphs. Clean blowup preserves log smoothness and normal-crossings/corners structures. The quoted form of local coordinates is the mixed analogue of the standard tree coordinates for Fulton-MacPherson spaces: every holomorphic difference factors into the product of scales along a holomorphic collision tree times a unit, and similarly every time difference factors through topological collision scales. The mixed blowups ensure compatibility when the same divergent subgraph controls both.
\end{proof}

\subsection{Logarithmic BPHZ for graph-valued forms}
For each graph $\Gamma$ let
\[
 \omega_{\Gamma,\varepsilon}
\]
be the UV-regularized graph form on $C_\Gamma(\mathbf z)$, valued in
\[
 \End(V_1\otimes\cdots\otimes V_n)\wot \mathrm B(E_n),
\]
constructed from regularized propagators in the sense of \cite{GKW25,DNP25}. Pulling back to a boundary chart of $\mathfrak C_\Gamma(\mathbf z)$, the form has an asymptotic expansion of the shape
\[
 \beta_F^*\omega_{\Gamma,\varepsilon}
 \sim
 \sum_{\alpha,\beta,\mu,\nu}
 r^{\alpha}s^{\beta}(\log r)^{\mu}(\log s)^{\nu}
 \eta_{\alpha,\beta,\mu,\nu}(\xi,\tau;\varepsilon),
\]
with smooth coefficient forms $\eta_{\alpha,\beta,\mu,\nu}$.

\begin{definition}[Logarithmic Taylor projector and log-BPHZ operator]\label{def:log-bphz}
For a forest face $F$, define the logarithmic Taylor projector $T_F$ by keeping the asymptotic terms with some negative scale exponent:
\[
 T_F(\omega)
 :=
 \sum_{\exists e:\,\alpha_e<0\ \text{or}\ \exists f:\,\beta_f<0}
 r^{\alpha}s^{\beta}(\log r)^{\mu}(\log s)^{\nu}
 \eta_{\alpha,\beta,\mu,\nu}.
\]
The \emph{log-BPHZ operator} is
\[
 R_\Gamma(\omega):=\sum_{F\,\mathrm{forest}}(-1)^{|F|}T_F(\omega).
\]
Provided the limit exists, the renormalized graph form is
\[
 \omega_\Gamma^{\mathrm{ren}}:=\lim_{\varepsilon\to 0}R_\Gamma(\omega_{\Gamma,\varepsilon}).
\]
\end{definition}

This is the logarithmic analogue of forest subtraction: one subtracts precisely the asymptotic pieces whose scale degree is too negative to be integrable.

\begin{theorem}[Existence, logarithmicity, and residue factorization; \ClaimStatusProvedHere]\label{thm:log-bphz-existence}
In the quasi-linear holomorphic-topological setting, for every connected admissible graph $\Gamma$ the renormalized graph form $\omega_\Gamma^{\mathrm{ren}}$ exists and extends to a log-polyhomogeneous differential form on $\mathfrak C_\Gamma(\mathbf z)$ with at worst simple logarithmic singularities:
\[
 \omega_\Gamma^{\mathrm{ren}}
 \in
 \Omega^{\bullet}(\mathfrak C_\Gamma(\mathbf z),\log\partial)
 \wot
 \End(V_1\otimes\cdots\otimes V_n)\wot \mathrm B(E_n).
\]
On a boundary divisor $D_F\subset\mathfrak C_\Gamma(\mathbf z)$ corresponding to a connected collision forest $F$,
\[
 \Res_{D_F}\,\omega_\Gamma^{\mathrm{ren}}
 =
 \mathrm{Comp}_F\Big(
 \omega_{\Gamma/F}^{\mathrm{ren}}
 \boxtimes
 \bboxprod_{H\in F}\omega_H^{\mathrm{ren}}
 \Big),
\]
where $\Gamma/F$ is the graph obtained by collapsing each cluster of $F$ and $\mathrm{Comp}_F$ is the induced bar/endomorphism composition map.
\end{theorem}

\begin{proof}
We proceed in three steps: existence of the renormalized limit via forest subtraction, logarithmicity of the extension via the normal Euler field, and the residue factorization formula via propagator splitting.

\medskip
\textbf{Step 1: Existence via forest subtraction (induction on loop order).}
We argue by strong induction on the loop order $L(\Gamma)$. The base case $L(\Gamma)=0$ (tree graphs) requires no renormalization: the unregularized form $\omega_\Gamma$ is already smooth on the interior, and every propagator edge contributes a Cauchy kernel $K(z) = dz/(z-w)$ (holomorphic, with a simple pole absorbed by the $d\log$ kernel) times a Heaviside factor $H(t) = \theta(t_i - t_j)$ (piecewise smooth in the topological direction). The product over edges of a tree is integrable on the compactified fibers because the total scaling degree of a tree with $|E(\Gamma)|$ edges at any collision cluster $S$ is
\[
 \delta_S(\Gamma) = |E_{\mathrm{int}}(S)| \cdot (-1) + |E_{\mathrm{int}}(S)| \cdot 1 = 0,
\]
where $E_{\mathrm{int}}(S)$ denotes the edges internal to the cluster $S$. The holomorphic Cauchy kernel contributes weight $-1$ in the holomorphic scale (a simple pole in $z_i - z_j$), and the differential $d(z_i - z_j)$ contributes weight $+1$. Thus the tree-level integrand has non-negative scaling degree at every collision stratum and extends to a smooth form on the compactification.

For the inductive step, suppose the theorem holds for all graphs of loop order strictly less than $L(\Gamma)$. A connected divergent subgraph $\gamma \subset \Gamma$ is a connected subgraph whose superficial degree of divergence $\delta(\gamma) < 0$. In the quasi-linear setting, each propagator $P_e$ has total scaling weight $0$ (the simple pole from the Cauchy kernel is compensated by the differential), so the superficial degree of divergence of a subgraph $\gamma$ with $L(\gamma)$ loops, $|V(\gamma)|$ internal vertices, and $|E(\gamma)|$ edges is
\[
 \delta(\gamma) = 2L(\gamma) - |E_{\mathrm{ext}}(\gamma)|,
\]
where $|E_{\mathrm{ext}}(\gamma)|$ is the number of external edges of $\gamma$, and the factor $2$ arises because each loop contributes a two-dimensional holomorphic integration. (In the topological direction, the ordered real integration contributes dimension $1$ per time variable but the Heaviside factors are piecewise constant and do not worsen convergence.) The divergent subgraphs are exactly those with $\delta(\gamma) < 0$, and they are nested or disjoint by the Zimmermann forest structure.

The logarithmic Taylor projector $T_F$ (Definition~\ref{def:log-bphz}) extracts from $\omega_{\Gamma,\varepsilon}$ the asymptotic terms of negative scaling degree in the scales $(r_e, s_f)$ associated to the forest $F$. The log-BPHZ operator $R_\Gamma = \sum_F (-1)^{|F|} T_F$ implements the forest subtraction by inclusion-exclusion over all forests of divergent connected subgraphs. By the standard BPHZ convergence theorem adapted to the polyhomogeneous setting, $R_\Gamma(\omega_{\Gamma,\varepsilon})$ has non-negative scaling degree at every collision stratum, hence the limit
\[
 \omega_\Gamma^{\mathrm{ren}} = \lim_{\varepsilon \to 0} R_\Gamma(\omega_{\Gamma,\varepsilon})
\]
exists. The subtraction terms $T_F(\omega_{\Gamma,\varepsilon})$ for any forest $F$ involve only the renormalized forms $\omega_\gamma^{\mathrm{ren}}$ for proper divergent subgraphs $\gamma \subsetneq \Gamma$ (which exist by the inductive hypothesis, since $L(\gamma) < L(\Gamma)$ for every proper subgraph) composed with the reduced graph $\Gamma/\gamma$. This closes the induction.

\medskip
\textbf{Step 2: Log-polyhomogeneity with simple logarithmic singularities.}
After the forest subtraction, the renormalized form $\omega_\Gamma^{\mathrm{ren}}$ has an asymptotic expansion near each boundary stratum of the form
\[
 \omega_\Gamma^{\mathrm{ren}} \sim \sum_{\substack{\alpha,\beta,\mu,\nu \\ \alpha_e \ge 0,\, \beta_f \ge 0\, \forall e,f}} r^\alpha s^\beta (\log r)^\mu (\log s)^\nu\, \eta_{\alpha,\beta,\mu,\nu}(\xi,\tau),
\]
with smooth coefficient forms $\eta_{\alpha,\beta,\mu,\nu}$. All scale exponents are non-negative by construction (the subtraction removed all negative-degree terms). We must show that the logarithmic powers are at most $\mu_e, \nu_f \le 1$ (simple logarithmic singularities).

This is where the specific structure of the quasi-linear holomorphic-topological propagator is decisive. Each propagator has the form $P_e = K(z_e) \cdot H(t_e)$ where $K(z_e) = d\log(z_e)$ is a logarithmic $1$-form and $H(t_e) = \theta(t_i - t_j)$ is a step function. A product of $|E(\Gamma)|$ such propagators contains at most $|E(\Gamma)|$ logarithmic $1$-forms. After pulling back to a boundary chart via the blowdown map $\beta_F$ and substituting $z_i - z_j = (\prod_{e \in P(i,j)} r_e) u_{ij}$, each $d\log(z_i - z_j)$ decomposes as
\[
 d\log(z_i - z_j) = \sum_{e \in P(i,j)} \frac{dr_e}{r_e} + d\log(u_{ij}).
\]
Thus the logarithmic $1$-forms $dr_e/r_e$ appear with multiplicity equal to the number of edges whose path passes through $e$. But each factor $dr_e/r_e$ is an exact logarithmic $1$-form on $\mathfrak C_\Gamma(\mathbf z)$ with a simple pole along $D_e = \{r_e = 0\}$. The forest subtraction $R_\Gamma$ removes negative powers of the scales but does not introduce new logarithmic powers: the subtraction terms $T_F$ are polynomial in $\log r_e$ and $\log s_f$ of the same order as the original asymptotic expansion, and the alternating sum cancels all terms of negative scale degree without increasing the logarithmic multiplicity.

To make this precise, apply the normal Euler field argument of Theorem~\ref{thm:local-residue-splitting-synthesis} (the planted-forest residue splitting theorem). In the notation of the mixed boundary chart, let $E = \sum_e r_e \partial_{r_e} + \sum_f s_f \partial_{s_f}$ be the total normal Euler field. On a monomial $r^\alpha s^\beta (\log r)^\mu (\log s)^\nu \eta$ with $|\alpha| + |\beta| > 0$, the Lie derivative satisfies $\mathcal{L}_E(\cdot) = (|\alpha| + |\beta|)(\cdot)$, and the contracting homotopy $h = (|\alpha|+|\beta|)^{-1}\iota_E$ satisfies $(dh + hd)\omega = \omega$. Thus positive-degree terms are exact and contractible, and the only non-contractible terms are those with $|\alpha| = |\beta| = 0$, i.e., the pure logarithmic residue terms. Since the propagators are sums of $dr_e/r_e$ (weight $0$ in $r_e$, but contributing a $d\log$ factor) and smooth terms, the asymptotic expansion at $r_e = 0$ has at most a single power of $\log r_e$ from each edge passing through the stratum. After the forest subtraction removes all strictly negative powers, what remains has non-negative integer powers of the scales decorated by at most simple logarithmic factors. Therefore $\omega_\Gamma^{\mathrm{ren}} \in \Omega^\bullet(\mathfrak C_\Gamma(\mathbf z), \log \partial)$.

\medskip
\textbf{Step 3: Residue factorization.}
Let $D_F$ be the boundary divisor corresponding to a connected collision forest $F$, with local scale coordinate $r_F$ (in the holomorphic direction; the topological case is analogous). We compute $\Res_{D_F}\,\omega_\Gamma^{\mathrm{ren}}$ by extracting the coefficient of $d\log r_F = dr_F/r_F$.

Decompose the edge set of $\Gamma$ as $E(\Gamma) = E_{\mathrm{int}}(F) \sqcup E_{\mathrm{ext}}(F)$, where $E_{\mathrm{int}}(F)$ consists of edges both of whose endpoints belong to the same cluster of $F$, and $E_{\mathrm{ext}}(F)$ consists of edges connecting distinct clusters (or connecting a cluster to an external vertex).

\emph{Internal edges.} For an internal edge $e \in E_{\mathrm{int}}(F)$ with endpoints in a cluster $H \in F$, the propagator $P_e = d\log(z_i - z_j) \cdot H(t_i - t_j)$ rescales entirely within the local coordinates of $H$. As $r_F \to 0$, the relative coordinates $\xi_{ij}$ within the cluster remain bounded (they parametrize the fiber $\FM_{|H|}^{\mathrm{red}}$), and the propagator restricts to the corresponding propagator of the subgraph $H$.

\emph{External edges.} For an external edge $e \in E_{\mathrm{ext}}(F)$ connecting cluster $H$ to an external point or a different cluster, the holomorphic difference $z_i - z_j$ factors as $r_F \cdot r_H \cdots u_{ij}$ where $r_H$ is the scale of the cluster containing $z_i$. As $r_F \to 0$, the external propagator freezes to its value at the collapsed configuration, contributing a smooth coefficient evaluated at the collision point of $H$.

The residue along $D_F$ therefore factorizes the graph $\Gamma$ into:
\begin{itemize}
 \item the renormalized subgraph forms $\omega_H^{\mathrm{ren}}$ for each cluster $H \in F$, integrated over the fiber $\FM_{|H|}^{\mathrm{red}}$ (the internal configuration within each cluster);
 \item the renormalized collapsed-graph form $\omega_{\Gamma/F}^{\mathrm{ren}}$ on the reduced configuration space, in which each cluster $H$ has been contracted to a single point.
\end{itemize}
The bar labels and endomorphism insertions compose via the contraction map $\mathrm{Comp}_F$: each cluster $H$ contributes a bar word that is composed with the bar word of the collapsed graph by the bar coalgebra structure, and the endomorphism factors compose by the tensor product structure of $\End(V_1 \otimes \cdots \otimes V_n)$. This composition is exactly the map $\mathrm{Comp}_F$ determined by the forest $F$.

It remains to verify that the forest subtraction is compatible with residue-taking: the residue of $R_\Gamma(\omega_{\Gamma,\varepsilon})$ along $D_F$ equals the composed renormalized forms of the sub- and quotient graphs. This follows from the Connes--Kreimer factorization property of the forest formula: for a forest $F' \neq F$, the subtraction $T_{F'}$ either refines $F$ (in which case it contributes a sub-forest subtraction within a cluster, giving the renormalized subgraph form by the inductive hypothesis) or is transverse to $F$ (in which case the subtraction commutes with residue extraction). The net effect is that $\Res_{D_F} \circ R_\Gamma = \mathrm{Comp}_F \circ (R_{\Gamma/F} \boxtimes \prod_{H \in F} R_H)$, which is the stated factorization formula.
\end{proof}

\subsection{Fiber integrals and the analytic collision field}
Let
\[
 \pi_\Gamma : \mathfrak C_\Gamma(\mathbf z) \longrightarrow X_n
\]
be the map forgetting all internal bulk and time variables. Define the graph weight
\[
 W_\Gamma := (\pi_\Gamma)_*\omega_\Gamma^{\mathrm{ren}}.
\]
Because the fibers are compactified and the integrands are logarithmic, the pushforward is again a logarithmic form on $X_n$ valued in the completed bar/endomorphism algebra.

Set
\[
 \Gamma_n^{\mathrm{an}} := \sum_{\Gamma\,\mathrm{connected}} \h^{L(\Gamma)}W_\Gamma,
\]
where $L(\Gamma)$ is the loop order of $\Gamma$. The sum is a formal $\h$-series. Define the analytic superconnection
\[
 \SN_n^{\mathrm{an}} := \ds+\widehat\Theta_n+b_n+I_{\Gamma_n^{\mathrm{an}}}
\]
on the trivial bar bundle over $X_n$.

\begin{theorem}[Analytic curvature formula; \ClaimStatusProvedHere]\label{thm:analytic-curvature}
The curvature of the analytic superconnection is
\[
 (\SN_n^{\mathrm{an}})^2 = I_{\mathcal K_n^{\mathrm{an}}},
\]
where
\[
 \mathcal K_n^{\mathrm{an}}
 =
 (\ds+\Theta_n)\Gamma_n^{\mathrm{an}}
 +
 \sum_{m\ge 1}m_m(\Gamma_n^{\mathrm{an}},\dots,\Gamma_n^{\mathrm{an}}).
\]
\end{theorem}

\begin{proof}
This is the formal master curvature formula of Theorem~\ref{thm:master-curvature} applied coefficientwise to the analytic collision field $\Gamma_n^{\mathrm{an}}$.
\end{proof}

\begin{theorem}[Compactified-Stokes flatness; \ClaimStatusProvedHere]\label{thm:analytic-flatness}
In the quasi-linear holomorphic-topological theory described above,
\[
 \mathcal K_n^{\mathrm{an}}=0.
\]
Equivalently,
\[
 (\SN_n^{\mathrm{an}})^2=0.
\]
\end{theorem}

\begin{proof}
By Theorem~\ref{thm:analytic-curvature}, the curvature $(\SN_n^{\mathrm{an}})^2 = I_{\mathcal K_n^{\mathrm{an}}}$ where
\[
 \mathcal K_n^{\mathrm{an}} = (\ds + \Theta_n)\Gamma_n^{\mathrm{an}} + \MC(\Gamma_n^{\mathrm{an}}).
\]
We must show that each graph-weight coefficient of $\mathcal K_n^{\mathrm{an}}$ vanishes. The proof combines three ingredients: Stokes' theorem on the compactified mixed configuration spaces $\mathfrak C_\Gamma(\mathbf z)$, the BV/BRST master equation in the interior, and the Arnold cancellation at codimension-two corners (the FM analogue of Lemma~\ref{lem:Arnold_cancellation}).

\medskip
\textbf{Step 1: Stokes decomposition for a single graph weight.}
Fix a connected admissible graph $\Gamma$ and consider the graph weight $W_\Gamma = (\pi_\Gamma)_* \omega_\Gamma^{\mathrm{ren}}$. The fiber integral is over the compact fibers of $\pi_\Gamma : \mathfrak C_\Gamma(\mathbf z) \to X_n$. By Theorem~\ref{thm:log-bphz-existence}, $\omega_\Gamma^{\mathrm{ren}}$ is a logarithmic form on $\mathfrak C_\Gamma(\mathbf z)$, so the Stokes formula (Theorem~\ref{thm:Stokes_FM}) applies:
\[
 (\pi_\Gamma)_* \bigl(d_{\mathrm{fiber}}\, \omega_\Gamma^{\mathrm{ren}}\bigr)
 =
 \sum_{D \subset \partial_{\mathrm{fiber}} \mathfrak C_\Gamma} \epsilon_D\, (\pi_\Gamma|_D)_* \Res_D\, \omega_\Gamma^{\mathrm{ren}},
\]
where the sum runs over all codimension-one boundary divisors of the fiber, $\epsilon_D$ is the orientation sign from the outward-normal-first convention, and $d_{\mathrm{fiber}}$ denotes the de Rham differential along the fiber.

The total exterior derivative of $W_\Gamma$ on the base $X_n$ satisfies
\begin{equation}\label{eq:stokes-graph-weight}
 (\ds + \Theta_n) W_\Gamma
 =
 (\pi_\Gamma)_*\bigl((d_{\mathrm{total}} + \Theta_n)\omega_\Gamma^{\mathrm{ren}}\bigr)
 +
 \sum_{D \subset \partial_{\mathrm{fiber}} \mathfrak C_\Gamma} \epsilon_D\, (\pi_\Gamma|_D)_* \Res_D\, \omega_\Gamma^{\mathrm{ren}},
\end{equation}
where $d_{\mathrm{total}} = d_{\mathrm{base}} + d_{\mathrm{fiber}}$ is the total de Rham differential and the boundary term arises from the Stokes contribution along the fiber boundary.

\medskip
\textbf{Step 2: Interior contribution and the perturbative master equation.}
The interior term $(\pi_\Gamma)_* ((d_{\mathrm{total}} + \Theta_n)\omega_\Gamma^{\mathrm{ren}})$ encodes the differential of the graph integrand in the bulk of the configuration space. In the perturbative BV/BRST framework organized by \cite{GKW25}, the graph-valued form $\omega_\Gamma^{\mathrm{ren}}$ is constructed from the propagator $P = K(z) \cdot H(t)$ and vertex insertions determined by the $A_\infty$ structure maps $m_k$. The BV master equation $Q_{\mathrm{BV}}^2 = 0$ for the theory on $\CC_z \times \RR_t$ implies that
\[
 (d_{\mathrm{total}} + \Theta_n + Q_{\mathrm{vertex}})\omega_\Gamma^{\mathrm{ren}}
 =
 \sum_{\gamma \subset \Gamma,\, |\gamma| \ge 2}
 \omega_{\Gamma/\gamma}^{\mathrm{ren}} \circ_\gamma \omega_\gamma^{\mathrm{ren}},
\]
where $Q_{\mathrm{vertex}}$ encodes the action of the bar differential $b_n$ on vertex labels, the sum runs over all connected subgraphs $\gamma$ that can be contracted, and $\circ_\gamma$ denotes the composition obtained by collapsing $\gamma$ and composing the resulting bar words. The key identity is that $Q_{\mathrm{BV}}^2 = 0$ holds on the interior of $\mathfrak C_\Gamma$ by the classical master equation of the BV theory, and renormalization preserves this identity because the forest subtraction is compatible with the BV differential (the subtraction terms are local counterterms annihilated by $Q_{\mathrm{BV}}$).

After pushforward, the interior contribution therefore produces terms of the form $\sum_\gamma W_{\Gamma/\gamma} \circ_\gamma W_\gamma$, which are components of the Maurer--Cartan polynomial $\MC(\Gamma_n^{\mathrm{an}})$.

\medskip
\textbf{Step 3: Boundary contribution and residue factorization.}
The boundary of the fiber $\partial_{\mathrm{fiber}} \mathfrak C_\Gamma$ consists of collision divisors $D_F$ indexed by connected collision forests $F$. By the residue factorization formula of Theorem~\ref{thm:log-bphz-existence},
\[
 \Res_{D_F}\, \omega_\Gamma^{\mathrm{ren}}
 =
 \mathrm{Comp}_F\Big(
 \omega_{\Gamma/F}^{\mathrm{ren}}
 \boxtimes
 \bboxprod_{H \in F} \omega_H^{\mathrm{ren}}
 \Big).
\]
After pushforward along $\pi_\Gamma|_{D_F}$, this yields
\[
 (\pi_\Gamma|_{D_F})_* \Res_{D_F}\, \omega_\Gamma^{\mathrm{ren}}
 =
 \mathrm{Comp}_F\Big(
 W_{\Gamma/F}
 \boxtimes
 \bboxprod_{H \in F} W_H
 \Big).
\]
Summing over all graphs $\Gamma$ and all collision forests $F$, the total boundary contribution reconstitutes the composition terms of the Maurer--Cartan polynomial: each composition $m_\ell(\Gamma_n^{\mathrm{an}}, \ldots, \Gamma_n^{\mathrm{an}})$ receives contributions from all graphs $\Gamma$ and all ways of extracting $\ell$ collision clusters from $\Gamma$. The signed sum of these contributions is exactly $\MC(\Gamma_n^{\mathrm{an}})$.

\medskip
\textbf{Step 4: Codimension-two corner cancellation.}
It remains to verify that the codimension-two contributions from corners where two collision divisors $D_F$ and $D_{F'}$ meet do not introduce extra terms. These corners arise in two cases:

\emph{Case (a): Disjoint collision forests.} If $F$ and $F'$ have disjoint vertex sets, the corner $D_F \cap D_{F'}$ parametrizes simultaneous but independent collapses. The two orderings of taking residues are related by transposing the logarithmic $1$-forms $d\log r_F$ and $d\log r_{F'}$. Since these are both degree-$1$ forms, the transposition contributes a sign $(-1)^{1 \cdot 1} = -1$:
\[
 \Res_{D_F}(\Res_{D_{F'}}(\omega_\Gamma^{\mathrm{ren}}))
 = -\Res_{D_{F'}}(\Res_{D_F}(\omega_\Gamma^{\mathrm{ren}})).
\]
The two contributions to the boundary integral at this corner appear with the same combined orientation sign $\epsilon_{D_F} \cdot \epsilon_{D_{F'}}$ (since the outward normals commute when the forests are disjoint), and they cancel:
\[
 \epsilon_{D_F}\, \Res_{D_F}(\Res_{D_{F'}}(\omega_\Gamma^{\mathrm{ren}}))
 + \epsilon_{D_{F'}}\, \Res_{D_{F'}}(\Res_{D_F}(\omega_\Gamma^{\mathrm{ren}}))
 = 0.
\]
This is the direct analogue of the Arnold cancellation at corners (Lemma~\ref{lem:Arnold_cancellation}) for the mixed logarithmic configuration spaces.

\emph{Case (b): Nested collision forests.} If $F \subset F'$ (one collision cluster is a refinement of another), the corner $D_F \cap D_{F'}$ is a codimension-$1$ boundary of $D_{F'}$, not a genuine codimension-$2$ corner of $\mathfrak C_\Gamma$. This contribution is already accounted for in the recursive structure: the residue along $D_{F'}$ produces a renormalized subgraph form, and the further collapse $F \subset F'$ is a boundary contribution \emph{within} that renormalized form. By the inductive structure of Theorem~\ref{thm:log-bphz-existence}, these nested contributions are consistently incorporated into the renormalized forms of the sub- and quotient graphs. No separate cancellation is required.

\medskip
\textbf{Step 5: Assembly.}
Combining Steps 1--4: the Stokes identity \eqref{eq:stokes-graph-weight} gives
\[
 0 = (\ds + \Theta_n)\Gamma_n^{\mathrm{an}} + \MC(\Gamma_n^{\mathrm{an}}) = \mathcal K_n^{\mathrm{an}}.
\]
The interior contributions (Step 2) produce the bar-differential terms and part of the Maurer--Cartan polynomial. The boundary contributions (Step 3) produce the remaining collision-composition terms, completing $\MC(\Gamma_n^{\mathrm{an}})$. The codimension-two corners cancel by the Arnold mechanism (Step 4). Since the total Stokes integral over each compact fiber vanishes ($\int_{\partial \mathrm{fiber}} = \int_{\mathrm{fiber}} d_{\mathrm{fiber}}$ and $d_{\mathrm{fiber}}^2 = 0$ ensure $\partial^2 = 0$ on the chain level), the assembled identity $\mathcal K_n^{\mathrm{an}} = 0$ follows.
\end{proof}

\begin{theorem}[Binary residues and the analytic $A_\infty$ Yang-Baxter identity; \ClaimStatusProvedHere]\label{thm:analytic-yb}
Let $D_{ij}\subset X_n$ be the boundary divisor where $z_i$ and $z_j$ collide. Then near $D_{ij}$,
\[
 \Gamma_n^{\mathrm{an}} = \mathbf r_{ij}(u)+\text{regular terms},
 \qquad u=z_i-z_j,
\]
where $\mathbf r_{ij}(u)$ is the analytic binary collision kernel obtained by integrating the codimension-one binary-collision faces.

On a codimension-two triple-collision corner one has the iterated-residue identity
\begin{align*}
 &\MC\big(\mathbf r_{12}(u)+\mathbf r_{13}(u+v)+\mathbf r_{23}(v)\big)
 -\MC\big(\mathbf r_{12}(u)\big)
 -\MC\big(\mathbf r_{13}(u+v)\big)
 -\MC\big(\mathbf r_{23}(v)\big)
 =0.
\end{align*}
\end{theorem}

\begin{proof}
The existence of the binary residue follows directly from the codimension-one part of Theorem~\ref{thm:log-bphz-existence}. For the codimension-two identity, apply Theorem~\ref{thm:analytic-flatness} and take the iterated residue along the triple-collision corner. Only the three compatible binary gluings survive, and their signed sum is exactly the displayed identity.
\end{proof}

\begin{theorem}[Family version over semistable degenerations; \ClaimStatusProvedHere]\label{thm:family-version}
Let $C\to B$ be a semistable degeneration of curves whose generic fiber is $\PP^1$. Then the construction above extends fiberwise to a logarithmic flat analytic superconnection on the relative logarithmic Fulton-MacPherson space $\FM_n(C/B)$. On a rigid special-fiber component indexed by $\rho$, the relative collision field factorizes after pullback along Mok's birational product model:
\[
 \Gamma_{n/B,\rho}^{\mathrm{an}}
 =
 \mathrm{Comp}_\rho\Big(
 \bboxprod_{v\in V(S_\rho)} \Gamma_{I_v}^{\mathrm{an}}(Y_v\mid D_v)
 \Big).
\]
\end{theorem}

\begin{proof}
The proof proceeds in four steps: relative log-smoothness of the compactified configuration space, extension of the graph-integral construction to the relative setting, logarithmic control of fiber integrals along the discriminant, and the factorization formula on rigid special-fiber components.

\medskip
\textbf{Step 1: Relative log-smoothness of $\FM_n(C/B)$.}
Let $\pi : C \to B$ be a semistable degeneration whose generic fiber is $\PP^1$. The total space $C$ is a smooth surface, $B$ is a smooth curve, and the special fibers are nodal curves $C_0 = Y_1 \cup_{p_1} Y_2 \cup_{p_2} \cdots \cup_{p_{m-1}} Y_m$ where each $Y_v \cong \PP^1$ and $p_j$ are the nodes. Equip $C$ with the divisorial log structure induced by the special fiber $C_0$ and equip $B$ with the log structure induced by the discriminant locus $\Delta = \{0\} \subset B$. Then $\pi : (C, M_C) \to (B, M_B)$ is a log-smooth morphism of fs log schemes.

The relative Fulton--MacPherson space $\FM_n(C/B)$ is constructed by Mok's procedure \cite{Mok25}: starting from the $n$-fold fiber product $C^n_B := C \times_B \cdots \times_B C$, one performs iterated clean blowups along the diagonals (including the relative diagonals and the loci where points approach the nodes), ordered by decreasing codimension. The result is a proper birational modification $\FM_n(C/B) \to C^n_B$ which is log smooth over $B$.

The log-smoothness is verified locally. Away from the nodes, $C/B$ is smooth, and $\FM_n(C/B)$ coincides with the standard relative Fulton--MacPherson space, which is smooth (and hence log smooth) by the classical construction. Near a node $p$ of the special fiber, choose local coordinates $(x, y, t)$ on $C$ such that $\pi(x,y,t) = t$ and $C_0 = \{xy = t\}$ locally. The $n$-fold fiber product has local equations $\{x_i y_i = t,\, i = 1, \ldots, n\}$ in $\AA^{2n+1}$. The relative diagonals $\{z_i = z_j\}$ are expressed in terms of the component coordinates: on the component $Y_v$, the identification $z_i = x_i$ or $z_i = y_i^{-1}$ gives well-defined relative positions. The iterated blowup along these loci and their intersections with the nodal strata produces a space with simple normal crossings boundary. In Mok's logarithmic coordinates, the boundary divisors of $\FM_n(C/B)$ are of two types:
\begin{itemize}
 \item \emph{collision divisors} $D_S$ indexed by subsets $S \subset [n]$ with $|S| \ge 2$, where points in $S$ collide on a smooth component of the fiber;
 \item \emph{degeneration divisors} $D_t = \pi^{-1}(\Delta)$, where the base parameter $t \to 0$ and the fiber degenerates.
\end{itemize}
These divisors have simple normal crossings, and $\FM_n(C/B) \to B$ is log smooth with respect to the divisorial log structures. Near a collision stratum within a smooth component, the local coordinates are exactly the FM scale coordinates $(r_e, \xi)$ of Proposition~\ref{prop:mixed-log-smooth}. Near the degeneration locus, there are additional scale coordinates $(\rho_j)_{j=1,\ldots,m-1}$ (the gluing parameters at the nodes), with $t = \prod_j \rho_j^{a_j}$ for appropriate multiplicities $a_j$, and unit coordinates parametrizing the relative positions on each component.

\medskip
\textbf{Step 2: Relative graph-integral construction.}
The mixed logarithmic configuration space $\mathfrak C_\Gamma(C/B)$ of Definition~\ref{def:mixed-log-config} extends to the relative setting: for each admissible graph $\Gamma$, the internal bulk vertices are placed on the fibers of $C/B$, and the mixed blowup construction is performed fiberwise, with further blowups at the loci where internal vertices approach the nodes. The result $\mathfrak C_\Gamma(C/B) \to B$ is a proper log-smooth family over $B$.

The regularized graph form $\omega_{\Gamma,\varepsilon}$ and the log-BPHZ operator $R_\Gamma$ of Definition~\ref{def:log-bphz} are defined fiberwise: for each $b \in B$, the restriction to the fiber $\mathfrak C_\Gamma(C_b)$ is constructed exactly as in the absolute case. The propagator $P_e = K(z_e) \cdot H(t_e)$ on each fiber extends to a logarithmic form on $\mathfrak C_\Gamma(C/B)$ because the Cauchy kernel $K(z) = d\log(z_i - z_j)$ is a relative logarithmic $1$-form (logarithmic along both the collision and degeneration divisors) and the Heaviside factor $H(t)$ depends only on the topological direction which is constant in the family.

The renormalized graph form $\omega_\Gamma^{\mathrm{ren}} = \lim_{\varepsilon \to 0} R_\Gamma(\omega_{\Gamma,\varepsilon})$ exists by the same forest-subtraction argument as Theorem~\ref{thm:log-bphz-existence}, applied fiberwise. The induction on loop order and the scaling-degree analysis are identical because the superficial degree of divergence is a local invariant that does not depend on the global geometry of the fiber.

\medskip
\textbf{Step 3: Logarithmic extension along the discriminant.}
The critical point is that the fiber integral
\[
 W_\Gamma(b) = (\pi_\Gamma)_*\omega_\Gamma^{\mathrm{ren}}\big|_b
\]
extends to a logarithmic form on $B$ with at worst simple logarithmic singularities along $\Delta$.

Near the discriminant $t = 0$, a fiber $C_t$ degenerates to the nodal curve $C_0 = \bigcup_v Y_v$. The relative FM space $\FM_n(C/B)$ has a boundary divisor $D_t = \{t = 0\}$ with local coordinate $t$. The renormalized graph form $\omega_\Gamma^{\mathrm{ren}}$ is logarithmic on $\mathfrak C_\Gamma(C/B)$ (by Step 2), so it has an expansion
\[
 \omega_\Gamma^{\mathrm{ren}} = \frac{dt}{t} \wedge \alpha + \beta
\]
with $\alpha, \beta$ smooth at $t = 0$. Pushing forward along the fibers of $\pi_\Gamma$, the fiber integral of $\omega_\Gamma^{\mathrm{ren}}$ decomposes as
\[
 W_\Gamma(b) = \frac{dt}{t} \cdot (\pi_\Gamma)_* \alpha\big|_{t=0} + (\pi_\Gamma)_* \beta + O(t \log t).
\]
The first term is a simple logarithmic singularity along $\Delta$. To verify that no worse singularities arise, note that the fiber $\mathfrak C_\Gamma(C_0)$ over $t = 0$ is a compact manifold with corners (a union of FM spaces on the components $Y_v$ glued along nodal strata), and the integrand $\alpha|_{t=0}$ is a smooth form on this compact fiber. Thus $(\pi_\Gamma)_* \alpha|_{t=0}$ is finite and smooth. By Proposition~\ref{prop:mixed-log-smooth}, the total space is log smooth, so the pushforward inherits the logarithmic structure of the base. Higher logarithmic powers $(\log t)^k$ with $k \ge 2$ are excluded by the same Euler-field argument as in Step 2 of Theorem~\ref{thm:log-bphz-existence}: the normal Euler field $t\, \partial_t$ provides a contracting homotopy for terms of positive normal degree, and the renormalization has already removed all negative-degree terms.

Therefore the total analytic collision field $\Gamma_{n/B}^{\mathrm{an}} = \sum_\Gamma \h^{L(\Gamma)} W_\Gamma$ is a formal $\h$-series of logarithmic forms on $(B, \Delta)$ valued in the completed bar/endomorphism algebra, and the analytic superconnection $\SN_{n/B}^{\mathrm{an}}$ is a well-defined logarithmic superconnection on the total base $B$.

Flatness of the relative superconnection, $(\SN_{n/B}^{\mathrm{an}})^2 = 0$, follows by applying the same argument as Theorem~\ref{thm:analytic-flatness} fiberwise: the Stokes identity, BV master equation, and Arnold corner cancellation hold on each fiber $\mathfrak C_\Gamma(C_b)$ (including at $b = 0$, where the fiber is a union of FM spaces on the components), and the resulting identity propagates to the total relative collision field by continuity of the fiber integral.

\medskip
\textbf{Step 4: Factorization on rigid special-fiber components.}
Let $C_0 = Y_1 \cup_{p_1} \cdots \cup_{p_{m-1}} Y_m$ be the special fiber, and let $\rho$ be a rigid special-fiber component, i.e., an assignment $I_v \subset [n]$ of marked points to components $Y_v$ (with $\bigsqcup_v I_v = [n]$) such that the configuration is stable on each component. Let $S_\rho$ denote the dual graph of $C_0$ decorated by the partition $(I_v)_v$. By Mok's birational product model \cite{Mok25}, the fiber of $\FM_n(C/B)$ over this rigid component admits a birational identification
\[
 \FM_n(C_0)\big|_\rho
 \cong
 \prod_{v \in V(S_\rho)} \FM_{I_v}(Y_v \mid D_v),
\]
where $D_v \subset Y_v$ is the set of nodal points on $Y_v$, and the birational modification resolves the loci where marked points on different components approach the same node.

Now consider a connected admissible graph $\Gamma$ with external vertices distributed according to the partition $(I_v)_v$. The graph decomposes as $\Gamma = \bigcup_v \Gamma_v \cup \Gamma_{\mathrm{node}}$, where $\Gamma_v$ consists of edges whose both endpoints lie on $Y_v$, and $\Gamma_{\mathrm{node}}$ consists of edges connecting different components (passing through nodes). As the degeneration parameter $t \to 0$:
\begin{itemize}
 \item The propagators for edges in $\Gamma_v$ restrict to the propagators on $(Y_v, D_v)$, since the holomorphic Cauchy kernel and topological Heaviside factor are local objects that depend only on the component.
 \item The propagators for edges in $\Gamma_{\mathrm{node}}$ degenerate: the holomorphic positions on different components become infinitely separated in the degeneration limit, and the corresponding Cauchy kernel $K(z_i - z_j) \to d\log(\rho_j) + \text{regular}$, where $\rho_j$ is the gluing parameter at the relevant node. This produces a residue contribution along the nodal divisor.
\end{itemize}

Taking the residue along the degeneration divisor and using the factorization property of Theorem~\ref{thm:log-bphz-existence} (applied to the collision forest consisting of the components of the special fiber), the relative collision field factorizes:
\[
 \Gamma_{n/B,\rho}^{\mathrm{an}}
 =
 \mathrm{Comp}_\rho\Big(
 \bboxprod_{v \in V(S_\rho)} \Gamma_{I_v}^{\mathrm{an}}(Y_v \mid D_v)
 \Big),
\]
where $\mathrm{Comp}_\rho$ is the composition map determined by the dual graph $S_\rho$: it composes the bar words from the individual components at the nodes, using the node-crossing propagator residues $\Res_{\rho_j = 0} K(z)$ as the gluing data. This composition is well-defined because the bar coalgebra structure is compatible with the componentwise decomposition (the bar coderivation $b_n$ respects the support decomposition indexed by the partition $(I_v)_v$), and the endomorphism factors compose by the tensor product structure.
\end{proof}

\section{Reduction to ordinary logarithmic monodromy}\label{sec:reduction}

The superconnection is the natural object before reduction, but the actual monodromic tensor structure lives on reduced line states. The central question is therefore: when does the bar-valued logarithmic superconnection collapse to an ordinary logarithmic connection on cohomology? The answer is: whenever the tree-level line complex is a genuine degree-zero resolution.

\subsection{Quantum concentration from classical concentration}
Assume from now on that the quantum line complex has the form
\[
 (E_n[[\h]],Q_n),
 \qquad
 Q_n=Q_{0,n}+\sum_{\ell\ge 1}\h^\ell Q_{n,\ell},
\]
where $E_n$ is a bounded finite-rank graded vector bundle over the relevant parameter space.

\begin{theorem}[Quantum concentration from classical concentration; \ClaimStatusProvedHere]\label{thm:quantum-concentration}
Suppose
\[
 H^j(E_n,Q_{0,n})=0 \qquad \text{for all } j\neq 0.
\]
Then
\[
 H^j(E_n[[\h]],Q_n)=0 \qquad \text{for all } j\neq 0.
\]
\end{theorem}

\begin{proof}
Filter $E_n[[\h]]$ by the complete decreasing filtration $F^p:=\h^pE_n[[\h]]$. Since $Q_n=Q_{0,n}+O(\h)$, the associated graded differential is $Q_{0,n}$. The resulting spectral sequence has
\[
 E_1^{p,q}\cong \h^p H^{p+q}(E_n,Q_{0,n}).
\]
By the hypothesis, $E_1^{p,q}$ vanishes unless $p+q=0$. Degree reasons force all higher differentials to vanish. Therefore the total cohomology is supported only in total degree $0$.
\end{proof}

\begin{corollary}[Flat $\h$-adic deformation of reduced states; \ClaimStatusProvedHere]\label{cor:flat-hbar}
If $H^0(E_n,Q_{0,n})$ has constant rank, then $H^0(E_n[[\h]],Q_n)$ is a flat $\h$-adic deformation of $H^0(E_n,Q_{0,n})$.
\end{corollary}

\subsection{Canonical $\h$-adic retract}

\begin{theorem}[Canonical $\h$-adic strong deformation retract; \ClaimStatusProvedHere]\label{thm:canonical-retract}
Assume $H^j(E_n,Q_{0,n})=0$ for $j\neq 0$ and $H^0(E_n,Q_{0,n})$ has constant rank. Then there is a canonical $\h$-adic strong deformation retract
\[
 (H_n,0) \underset{p_n}{\stackrel{i_n}{\rightleftarrows}} (E_n[[\h]],Q_n),
 \qquad
 Q_n h_n+h_nQ_n = \id - i_np_n,
\]
where $H_n:=H^0(E_n[[\h]],Q_n)$.
\end{theorem}

\begin{proof}
Choose a Hermitian metric on $E_n$. Let $Q_0:=Q_{0,n}$ and
\[
 \Delta_0 := Q_0Q_0^{\dagger}+Q_0^{\dagger}Q_0.
\]
Because $H^0(E_n,Q_0)$ has constant rank, the harmonic bundle
\[
 H_{0,n}:=\ker\Delta_0\subset E_n^0
\]
is a smooth subbundle. Let $p_{0,n}$ denote orthogonal projection onto $H_{0,n}$, $i_{0,n}$ the inclusion, and $G_{0,n}$ the Moore-Penrose inverse of $\Delta_0$ on $(\ker\Delta_0)^\perp$. Set
\[
 h_{0,n}:=Q_0^{\dagger}G_{0,n}.
\]
Then
\[
 Q_0h_{0,n}+h_{0,n}Q_0 = \id - i_{0,n}p_{0,n}.
\]

Now write $Q_n=Q_0+\Delta_n$ with $\Delta_n\in \h\,\End^1(E_n)[[\h]]$. Since $\Delta_n$ raises $\h$-adic order, the operator $1-h_{0,n}\Delta_n$ is invertible in the $\h$-adic topology. Define
\[
 \Sigma_n := \Delta_n(1-h_{0,n}\Delta_n)^{-1}.
\]
The basic perturbation lemma gives
\[
 i_n=i_{0,n}+h_{0,n}\Sigma_n i_{0,n},
 \qquad
 p_n=p_{0,n}+p_{0,n}\Sigma_n h_{0,n},
 \qquad
 h_n=h_{0,n}+h_{0,n}\Sigma_n h_{0,n},
\]
which define a strong deformation retract of $(E_n[[\h]],Q_n)$ onto $(H_{0,n}[[\h]],q_n)$ with transferred differential
\[
 q_n:=p_{0,n}\Sigma_n i_{0,n}.
\]
But $q_n$ has cohomological degree $+1$, while $H_{0,n}$ lies in degree $0$, so $q_n=0$. Thus the target complex is $(H_{0,n}[[\h]],0)$, which identifies canonically with $H_n$ by Theorem~\ref{thm:quantum-concentration}.
\end{proof}

\begin{proposition}[Analytic convergence under norm smallness; \ClaimStatusProvedHere]\label{prop:analytic-retract}
On any open set on which the exact-resummed perturbation satisfies
\[
 \|h_{0,n}\Delta_n\|<1
\]
on every finite-rank truncation, the formulas of Theorem~\ref{thm:canonical-retract} converge as honest analytic maps and not only as formal $\h$-adic series.
\end{proposition}

\begin{proof}
The $\h$-adic Neumann series for $(1-h_{0,n}\Delta_n)^{-1}$ converges absolutely under the stated operator-norm bound, and the perturbation formulas are finite compositions of that inverse with bounded maps.
\end{proof}

\subsection{One-form rigidity after reduction}
Let the analytic superconnection on $E_n[[\h]]$ be written as
\[
 \SN_n = \ds-\big(Q_n+\alpha_n\big),
 \qquad
 \alpha_n:=\sum_{k\ge 1}A_n^{(k)},
\]
with
\[
 A_n^{(k)}\in \Omega^k(X_n,\log\partial X_n)\ot \End^{1-k}(E_n)[[\h]].
\]
Transfer this superconnection along the canonical retract.

\begin{theorem}[One-form rigidity; \ClaimStatusProvedHere]\label{thm:one-form-rigidity}
The transferred superconnection on $H_n$ is an ordinary logarithmic flat connection,
\[
 \SN_n^{\mathrm{red}} = \ds-\A_n,
\]
with
\[
 \A_n = p_nA_n^{(1)}i_n.
\]
In particular, every higher-form piece $A_n^{(k)}$ with $k\ge 2$ disappears directly after reduction.
\end{theorem}

\begin{proof}
The standard homological perturbation formula for the transferred superconnection is
\[
 \SN_n^{\mathrm{red}}
 =
 \ds-\sum_{m\ge 1} p_n\alpha_n(h_n\alpha_n)^{m-1}i_n.
\]
Consider one summand
\[
 p_n A_n^{(k_1)} h_n A_n^{(k_2)} h_n \cdots h_n A_n^{(k_m)} i_n.
\]
Its form degree is $r=k_1+\cdots+k_m$, while its internal degree is
\[
 (1-k_1)+\cdots+(1-k_m)-(m-1)=1-r.
\]
An endomorphism of $H_n$, which is concentrated in internal degree $0$, must have internal degree $0$. Therefore $1-r=0$, so $r=1$. Since each $k_j\ge 1$, the only possibility is $m=1$ and $k_1=1$. Hence the only surviving term is
\[
 \A_n = p_nA_n^{(1)}i_n.
\]
The transferred superconnection is flat because the original one is flat. Since it has only a $1$-form part, it is an ordinary flat logarithmic connection.
\end{proof}

\begin{corollary}[Reduced residues; \ClaimStatusProvedHere]\label{cor:reduced-residues}
For every boundary divisor $D\subset X_n$,
\[
 \Res_D(\A_n)=p_n\Res_D(A_n^{(1)})i_n.
\]
\end{corollary}

\begin{proof}
Residues commute with projection, inclusion, and restriction because these maps are logarithmically regular.
\end{proof}

\subsection{Tangential basepoints and the HT associator}
Write near a normal-crossings corner of $X_n$
\[
 \A_n = \sum_{a=1}^r \Omega_a\,d\log t_a + \A_n^{\mathrm{reg}},
\]
where $t_1,\dots,t_r$ are local boundary coordinates and $\A_n^{\mathrm{reg}}$ is regular.

\begin{lemma}[Commuting corner residues; \ClaimStatusProvedHere]\label{lem:commuting-residues}
At every normal-crossings corner one has
\[
 [\Omega_a,\Omega_b]=0\qquad (1\le a,b\le r).
\]
\end{lemma}

\begin{proof}
The coefficient of $d\log t_a\wedge d\log t_b$ in the curvature of the flat connection $\ds-\A_n$ is exactly $[\Omega_a,\Omega_b]$, so it must vanish.
\end{proof}

\begin{lemma}[Tangential regularization; \ClaimStatusProvedHere]\label{lem:tangential-regularization}
For a tangential direction $v$ at a boundary corner, the gauge transform
\[
 G := \prod_{a=1}^r t_a^{-\Omega_a}
\]
regularizes the connection, and the resulting regularized transport to the tangential basepoint is finite and coordinate independent up to canonical conjugacy.
\end{lemma}

\begin{proof}
By Lemma~\ref{lem:commuting-residues}, the factors $t_a^{-\Omega_a}$ commute, so $G$ is well defined. A direct computation gives
\[
 G^{-1}(\ds-\A_n)G = \ds-\A_n^{\mathrm{reg}},
\]
which is regular at the corner. Parallel transport for the regular gauge therefore has a finite limit along any tangent ray representing $v$. Changing logarithmic coordinates multiplies $G$ by a regular invertible factor with a finite boundary value, so the regularized transport changes only by conjugation with that finite value.
\end{proof}

Fix three reduced line states $X,Y,Z$. On $X_3$ there are two canonical tangential basepoints:
\begin{itemize}
 \item $b_L$, the zone $|z_1-z_2|\ll |z_2-z_3|$, corresponding to $(XY)Z$;
 \item $b_R$, the zone $|z_2-z_3|\ll |z_1-z_3|$, corresponding to $X(YZ)$.
\end{itemize}

\begin{definition}[HT associator and braiding]\label{def:ht-assoc}
The \emph{HT associator} is the regularized parallel transport
\[
 \Phi_{\HT}(X,Y,Z):=\PT^{\mathrm{reg}}_{b_L\to b_R}(\A_3).
\]
For two reduced line states $X,Y$, the \emph{HT braiding} $c_{X,Y}$ is the half-monodromy of $\A_2$ along a counterclockwise exchange path in $\Conf_2(\CC)$.
\end{definition}

\begin{theorem}[Pentagon identity; \ClaimStatusProvedHere]\label{thm:pentagon}
The operators $a_{X,Y,Z}:=\Phi_{\HT}(X,Y,Z)$ satisfy the pentagon relation
\[
 (\id_X\ot a_{Y,Z,W})\,a_{X,Y\ot Z,W}\,(a_{X,Y,Z}\ot \id_W)
 =
 a_{X,Y,Z\ot W}\,a_{X\ot Y,Z,W}.
\]
\end{theorem}

\begin{proof}
Fix four ordered real points $0<z_2<z_3<1$, with $z_1=0$ and $z_4=1$. Their compactified real collision locus inside $X_4$ is a pentagon $K_4$: its five vertices correspond to the five parenthesizations of four letters, and its edges correspond to elementary reassociations. Restrict the reduced flat logarithmic connection to the interior of $K_4$. By Lemma~\ref{lem:tangential-regularization}, each edge carries a well-defined regularized transport between its tangential endpoints. The boundary of the pentagon is contractible, so the product of the five edge transports is the identity. Reading those edge transports in the standard cyclic order yields exactly the pentagon relation.
\end{proof}

\begin{theorem}[Hexagon identities; \ClaimStatusProvedHere]\label{thm:hexagon}
The braidings $c_{X,Y}$ and associators $a_{X,Y,Z}$ satisfy the two hexagon identities. One of them is
\[
 a_{Y,Z,X}\,c_{X,Y\ot Z}\,a_{X,Y,Z}
 =
 (\id_Y\ot c_{X,Z})\,a_{Y,X,Z}\,(c_{X,Y}\ot \id_Z).
\]
\end{theorem}

\begin{proof}
Both sides are regularized transports for the same reduced flat connection on the three-point configuration space. The left side moves $X$ around the fused cluster $Y\ot Z$; the right side first reassociates, then moves $X$ successively around $Y$ and $Z$, and finally reassociates back. These two composite motions are homotopic relative to their tangential endpoints. Flatness implies equality of their transports. The second hexagon is obtained by the mirror homotopy.
\end{proof}

\begin{corollary}[Monodromic braided tensor structure; \ClaimStatusProvedHere]\label{cor:braided-category}
The reduced line operators form a monodromic braided monoidal category: the tensor product is the reduced OPE, the associator is $\Phi_{\HT}$, and the braiding is given by half-monodromy.
\end{corollary}

\begin{theorem}[Local monodromy and pure braid representations; \ClaimStatusProvedHere]\label{thm:pure-braid}
For each binary collision divisor $D_{ij}\subset X_n$, let
\[
 \Omega_{ij}^{\mathrm{red}}:=\Res_{D_{ij}}(\A_n).
\]
The local monodromy around a small loop encircling $D_{ij}$ is conjugate to
\[
 \exp(2\pi i\,\Omega_{ij}^{\mathrm{red}}).
\]
Consequently the reduced connection defines a pure braid group representation
\[
 \rho_n: \pi_1(\Conf_n(\CC)) \longrightarrow \Aut(H_n).
\]
\end{theorem}

\begin{proof}
Near $D_{ij}$ choose a local coordinate $t$ such that
\[
 \A_n=\Omega_{ij}^{\mathrm{red}}\,\frac{dt}{t}+\text{regular terms}.
\]
After a local logarithmic gauge transformation removing the regular part, the monodromy around $t=\varepsilon e^{2\pi is}$ is exactly $\exp(2\pi i\,\Omega_{ij}^{\mathrm{red}})$. Since the connection is flat, parallel transport depends only on homotopy class and therefore defines a representation of the fundamental group.
\end{proof}

\begin{corollary}[Logarithmic periods; \ClaimStatusProvedHere]\label{cor:periods}
If the reduced connection and tangential basepoints are defined over a subfield $K\subset\CC$, then every matrix coefficient of $\Phi_{\HT}$ is a $K$-linear combination of regularized Chen iterated integrals of logarithmic $1$-forms on $X_3\setminus\partial X_3$.
\end{corollary}

\begin{proof}
Expand the path-ordered exponential for the regularized transport of the logarithmic connection. Each coefficient is an iterated integral of logarithmic one-forms between tangential basepoints, with coefficients in $K$.
\end{proof}

\section{Resolved degree-zero line systems and unconditionality}\label{sec:resolved}

The previous section reduces the monodromic theory to one concrete classical question: which physically natural line operators admit tree-level BRST complexes that are genuine degree-zero resolutions? Once that classical algebra is in place, the quantum-logarithmic monodromy machine is forced.

\begin{definition}[Resolved degree-zero line system]\label{def:resolved-line}
A tree-level line operator system is called \emph{resolved in degree $0$} if its BRST complex $(E_n,Q_{0,n})$ is a bounded finite-rank complex with
\[
 H^j(E_n,Q_{0,n})=0 \qquad (j\neq 0).
\]
\end{definition}

\subsection{Regular-sequence Koszul models}
Let $R$ be a Noetherian commutative ring, let $U$ be finite dimensional, and consider the Koszul complex
\[
 E = K_R(f_1,\dots,f_r)\ot U
 = \big(R\ot \Lambda(\eta_1,\dots,\eta_r),\;Q_0=\sum_{a=1}^r f_a\frac{\partial}{\partial\eta_a}\big)\ot U.
\]

\begin{proposition}[Regular-sequence resolution criterion; \ClaimStatusProvedHere]\label{prop:regular-sequence}
If $f_1,\dots,f_r$ is a regular sequence in $R$, then
\[
 H^j(E,Q_0)=0\quad (j\neq 0),
 \qquad
 H^0(E,Q_0)\cong (R/(f_1,\dots,f_r))\ot U.
\]
In particular, such a line system is resolved in degree $0$.
\end{proposition}

\begin{proof}
This is the standard exactness property of the Koszul complex associated to a regular sequence.
\end{proof}

\begin{lemma}[Explicit linear homotopy; \ClaimStatusProvedHere]\label{lem:linear-homotopy}
If $R=\kk[x_1,\dots,x_r]$, $f_a=x_a$, and
\[
 Q_0=\sum_{a=1}^r x_a\frac{\partial}{\partial\eta_a},
\]
then on positive total polynomial-plus-exterior degree the operator
\[
 h_0 := N^{-1}\sum_{a=1}^r \eta_a\frac{\partial}{\partial x_a},
 \qquad
 N:=\sum_{a=1}^r\Big(x_a\frac{\partial}{\partial x_a}+\eta_a\frac{\partial}{\partial\eta_a}\Big),
\]
satisfies
\[
 Q_0h_0+h_0Q_0=\id-\pi_0,
\]
where $\pi_0$ projects to degree $0$.
\end{lemma}

\begin{proof}
Set
\[
 \delta:=\sum_{a=1}^r \eta_a\frac{\partial}{\partial x_a}.
\]
A direct computation gives $[Q_0,\delta]=N$. Since $N$ is invertible on positive total degree, the operator $h_0=N^{-1}\delta$ satisfies the claim.
\end{proof}

\begin{example}[Complete-intersection line constraints]
Suppose a tree-level line coupling imposes equations $f_1=\cdots=f_r=0$ on an affine target, with $(f_1,\dots,f_r)$ a regular complete-intersection ideal. The associated BRST complex is a Koszul complex of the type above, hence is resolved in degree $0$. Such lines are therefore automatic inputs for the reduced monodromic theory.
\end{example}

\subsection{Unconditionality in the quasi-linear class}

\begin{theorem}[Unconditionality for resolved quasi-linear lines; \ClaimStatusProvedHere]\label{thm:unconditionality}
Consider a quasi-linear interacting holomorphic-topological theory in the sense of Section~\ref{sec:analytic}. Let $L_1,\dots,L_n$ be perturbative line operators whose tree-level BRST complexes are resolved in degree $0$. Then:
\begin{enumerate}[label=\textup{(\roman*)}]
 \item the full quantum line complexes remain concentrated in degree $0$;
 \item the canonical $\h$-adic retract exists;
 \item the reduced logarithmic flat connection is
 \[
 \A_n = p_nA_n^{(1)}i_n;
 \]
 \item the HT associator, braiding, pentagon, hexagon, and pure braid representations exist without any further concentration or retract hypothesis.
\end{enumerate}
\end{theorem}

\begin{proof}
Item (i) is Theorem~\ref{thm:quantum-concentration}. Item (ii) is Theorem~\ref{thm:canonical-retract}. Item (iii) is Theorem~\ref{thm:one-form-rigidity}. Item (iv) then follows from the transport constructions and coherence theorems of Section~\ref{sec:reduction}.
\end{proof}

\begin{theorem}[Boundary factorization after reduction; \ClaimStatusProvedHere]\label{thm:reduced-boundary-factorization}
Let $\rho$ be a rigid boundary component in a logarithmic degeneration of configuration spaces, with birational product model
\[
 X_n(\rho) \dashrightarrow \prod_{v\in V(S_\rho)} X_{I_v}.
\]
Assume the raw logarithmic $1$-form factorizes on that component:
\[
 A_n^{(1)}\big|_{X_n(\rho)}
 =
 \mathrm{Comp}_\rho\Big(\bboxprod_{v\in V(S_\rho)} A_{I_v}^{(1)}\Big).
\]
Then the reduced connection factorizes as
\[
 \A_n\big|_{X_n(\rho)}
 =
 \mathrm{Comp}_\rho\Big(\bboxprod_{v\in V(S_\rho)} \A_{I_v}\Big).
\]
\end{theorem}

\begin{proof}
On a disjoint cluster decomposition, the tree-level differential splits as a sum of the cluster differentials, so the Hodge projectors and inclusions tensorize:
\[
 p_n=\otimes_v p_{I_v},
 \qquad
 i_n=\otimes_v i_{I_v}.
\]
By one-form rigidity,
\[
 \A_n=p_nA_n^{(1)}i_n.
\]
Substituting the factorization of $A_n^{(1)}$ and tensorized projectors/inclusions gives the formula.
\end{proof}

\section{Arithmetic, enumerative, periodic, and chromatic consequences}\label{sec:consequences}

This section records what the theory says, in its own natural language, about the four directions that motivated the thread: arithmetic geometry, enumerative geometry, periodicity, and chromaticity.

\subsection{Arithmetic geometry: logarithmic local systems and periods}
The most immediate arithmetic output is that the reduced theory naturally lives in the category of logarithmic local systems on logarithmic Fulton-MacPherson spaces.

\begin{proposition}[Arithmetic form of the reduced theory; \ClaimStatusProvedHere]\label{prop:arithmetic-form}
Suppose the quasi-linear HT datum, the resolved line complexes, and the chosen retracts are all defined over a subfield $K\subset\CC$. Then for each $n$ the reduced logarithmic flat connection
\[
 \ds-\A_n
\]
on $X_n$ is defined over $K$, its local monodromies are exponentials of residue operators defined over $K$, and the HT associator takes values in the algebra of logarithmic periods over $K$ generated by iterated integrals of the entries of $\A_3$.
\end{proposition}

\begin{proof}
Everything is functorial in the coefficients. The residue statement is Corollary~\ref{cor:reduced-residues}, and the period statement is Corollary~\ref{cor:periods}.
\end{proof}

\subsection{Enumerative geometry: decorated collision trees}
On the strict side, Proposition~\ref{prop:strict-log-ext} attaches residue operators to the internal edges of collision trees. On the analytic side, Theorem~\ref{thm:log-bphz-existence} and Theorem~\ref{thm:family-version} say that degeneration strata are weighted by renormalized graph forms and factorization maps.

\begin{proposition}[Decorated degeneration formula; \ClaimStatusProvedHere]\label{prop:decorated-degeneration}
Every rigid boundary component of a logarithmic degeneration of configuration spaces carries a canonical operator-valued decoration obtained by composing the residue and factorization maps of the lower-point pieces. Thus the special fiber is not merely a sum over combinatorial trees; it is a sum over collision trees decorated by Yangian and $A_\infty$ residue data.
\end{proposition}

\begin{proof}
Combine the residue factorization theorem (Theorem~\ref{thm:log-bphz-existence}) with the family factorization theorem (Theorem~\ref{thm:family-version}) and the reduced boundary factorization theorem (Theorem~\ref{thm:reduced-boundary-factorization}).
\end{proof}

\subsection{Spectralization, periodicity, and chromaticity}
The classical chiral/factorization theory is linear over chain complexes and does not by itself impose a chromatic filtration. The correct chromatic object is the spectral lift.

\begin{theorem}[Smashing localization commutes with factorization homology; \ClaimStatusProvedHere]\label{thm:smashing-localization}
Let $A$ be a locally constant spectral factorization algebra on an $m$-manifold $M$, equivalently a $\Disk_m$-algebra in spectra. Let $L$ be a smashing localization of spectra. Then there is a canonical equivalence
\[
 L\Big(\int_M A\Big) \simeq \int_M LA.
\]
\end{theorem}

\begin{proof}
Factorization homology is computed as a colimit over disk covers:
\[
 \int_M A \simeq \operatorname*{colim}_{U\in \Disk(M)} A(U).
\]
If $L$ is smashing, then $L(X)\simeq E\ot X$ for some spectrum $E$, so $L$ preserves all colimits and is symmetric monoidal. Therefore
\[
 L\Big(\int_M A\Big)
 \simeq
 L\Big(\operatorname*{colim}_{U\in \Disk(M)}A(U)\Big)
 \simeq
 \operatorname*{colim}_{U\in \Disk(M)}LA(U)
 \simeq
 \int_M LA.
\]
\end{proof}

\begin{corollary}[Periodicity is local-to-global stable; \ClaimStatusProvedHere]\label{cor:periodicity}
If $LA$ takes values in $d$-periodic $L$-local spectra, then $\int_M LA$ is $d$-periodic.
\end{corollary}

\begin{proof}
Every object in the diagram computing $\int_M LA$ is $d$-periodic, and colimits in the $L$-local category preserve that periodicity.
\end{proof}

\begin{remark}[Chromaticity]
Theorem~\ref{thm:smashing-localization} says that chromatic height, once introduced by spectralization, is compatible with the local-to-global passage: chromaticity is not visible in the classical linear theory, but it appears after passing to spectral coefficients, where it behaves cleanly under factorization homology.
\end{remark}

\section{The Platonic ideal: resolved logarithmic HT data}\label{sec:platonic}

The constructions of the previous sections axiomatize independently of their perturbative realization.

\begin{definition}[Resolved logarithmic HT datum]\label{def:rlht-datum}
A \emph{resolved logarithmic holomorphic-topological datum} consists of the following.
\begin{enumerate}[label=\textup{(\roman*)}]
 \item For each $n\ge 1$, a logarithmic configuration space $X_n$ equipped with symmetric-group equivariance, boundary divisors indexed by collision types, and semistable family versions over degenerations.
 \item For each ordered $n$-tuple of line labels, a bounded graded vector bundle $E_n$ on $X_n$ with a differential $Q_n$ and a flat logarithmic superconnection
 \[
 \SN_n = \ds-(Q_n+\alpha_n),
 \]
 satisfying supportwise or Stokeswise factorization along boundary strata.
 \item A tree-level degree-zero resolution condition on $(E_n,Q_{0,n})$, or equivalently a canonical reduction datum
 \[
 (H_n,0) \underset{p_n}{\stackrel{i_n}{\rightleftarrows}} (E_n[[\h]],Q_n)
 \]
 compatible with boundary factorization.
 \item A translation structure and binary residues compatible with the collision divisors.
 \item A spectral lift to factorization algebras in spectra whenever periodicity or chromaticity is desired.
\end{enumerate}
\end{definition}

\begin{theorem}[Platonic ideal theorem; \ClaimStatusProvedHere]\label{thm:platonic-ideal}
Every resolved logarithmic HT datum determines canonically:
\begin{enumerate}[label=\textup{(\roman*)}]
 \item reduced logarithmic flat connections $\ds-\A_n$ on $H_n$;
 \item regularized tangential transports and hence associators and braidings;
 \item a monodromic braided monoidal category of reduced line states;
 \item pure braid group representations on reduced tensor powers;
 \item logarithmic period torsors over the field of definition;
 \item boundary factorization formulas on rigid components of semistable degenerations;
 \item after spectralization, periodic and chromatic local-to-global compatibilities.
\end{enumerate}
The concrete quasi-linear perturbative HT construction of Sections~\ref{sec:analytic}--\ref{sec:resolved} provides an example of such a datum whenever the tree-level line operators are resolved in degree $0$.
\end{theorem}

\begin{proof}
Items (i)--(iv) are exactly the contents of Theorems~\ref{thm:one-form-rigidity}, \ref{thm:pentagon}, \ref{thm:hexagon}, and \ref{thm:pure-braid}. Item (v) follows from Corollary~\ref{cor:periods}. Item (vi) follows from Theorems~\ref{thm:family-version} and \ref{thm:reduced-boundary-factorization}. Item (vii) follows from Theorem~\ref{thm:smashing-localization} and Corollary~\ref{cor:periodicity}. The final sentence follows from Theorem~\ref{thm:unconditionality}.
\end{proof}

\begin{remark}[The sharpened frontier question]
The unresolved frontier is now sharp: not whether a retract exists after quantization, but which physically natural line operators admit classically resolved degree-zero BRST models. Once that classical algebra exists, the quantum-logarithmic monodromy machine is forced.
\end{remark}

\section{Supplementary computations}

\subsection{Support bookkeeping in the Maurer-Cartan polynomial}
The support bookkeeping behind Theorem~\ref{thm:support-hierarchy} is as follows. If each $D_S$ is supported on the tensor factors in $S$, then the support of a higher product
\[
 m_\ell(D_{S_1},\dots,D_{S_\ell})
\]
is the union $S_1\cup\cdots\cup S_\ell$. Therefore the support-$S$ component of
\[
 \MC(\Gamma_n)=\sum_{\ell\ge 1}m_\ell(\Gamma_n,\dots,\Gamma_n)
\]
is exactly the sum over all tuples of nonempty subsets whose union is $S$.

\subsection{Degree count in one-form rigidity}
The proof of Theorem~\ref{thm:one-form-rigidity} uses only two facts:
\begin{enumerate}[label=\textup{(\alph*)}]
 \item $A_n^{(k)}$ has form degree $k$ and internal degree $1-k$;
 \item the homotopy $h_n$ has internal degree $-1$ and form degree $0$.
\end{enumerate}
Hence a term with $m$ insertions and total form degree $r$ has internal degree $1-r$. On degree-zero cohomology it can survive only if $r=1$, which forces a single $1$-form insertion and excludes all higher-form contributions.

\subsection{Why the codimension-two triple residue is Yang-Baxter}
At a triple collision corner there are exactly three codimension-one approaches to the corner, corresponding to the binary clusterings $(12)3$, $(13)2$, and $1(23)$. The compactified-Stokes boundary identity for the codimension-two corner is therefore the signed sum of the three corresponding binary gluing maps. Writing those three gluing maps in terms of the binary residue kernels gives the analytic $A_\infty$ Yang-Baxter identity of Theorem~\ref{thm:analytic-yb}.

\subsection{Explicit monodromy for $\widehat{\mathfrak{sl}}_2$: from genus zero to genus one}
\label{subsec:explicit-monodromy-sl2}

We carry out the full monodromy computation for the $\mathfrak{sl}_2$ KZ connection
at general level~$k$, starting from the Casimir, passing through
$\Conf_2$ and $\Conf_3$, extending logarithmically to $\FM_2(\PP^1\mid\infty)$,
and lifting to genus~$1$ via the Knizhnik--Zamolodchikov--Bernard (KZB) connection.
Throughout, $V=\CC^2$ denotes the fundamental representation of $\mathfrak{sl}_2$
with standard basis $v_+,v_-$ and generators
\[
 e=\begin{pmatrix}0&1\\0&0\end{pmatrix},\qquad
 f=\begin{pmatrix}0&0\\1&0\end{pmatrix},\qquad
 h=\begin{pmatrix}1&0\\0&-1\end{pmatrix}.
\]
The dual Coxeter number is $h^\vee=2$.

\begin{computation}[Casimir spectrum; \ClaimStatusProvedHere]
\label{comp:sl2-casimir-spectrum}
\index{Casimir!sl2@$\mathfrak{sl}_2$}
The split Casimir (invariant tensor) of $\mathfrak{sl}_2$ in the
Cartan--Weyl basis is
\[
 \Omega \;=\; e\otimes f + f\otimes e + \tfrac{1}{2}\,h\otimes h
 \;\in\; \mathfrak{sl}_2\otimes\mathfrak{sl}_2.
\]
On $V\otimes V$, using the ordered basis
$(v_+\!\otimes\!v_+,\;v_+\!\otimes\!v_-,\;v_-\!\otimes\!v_+,\;v_-\!\otimes\!v_-)$,
$\Omega$ is represented by
\[
 \Omega_{12} \;=\;
 \begin{pmatrix}
 \frac{1}{2} & 0 & 0 & 0 \\[3pt]
 0 & -\frac{1}{2} & 1 & 0 \\[3pt]
 0 & 1 & -\frac{1}{2} & 0 \\[3pt]
 0 & 0 & 0 & \frac{1}{2}
 \end{pmatrix}.
\]
The Clebsch--Gordan decomposition $V\otimes V\cong V_1\oplus V_0$ (triplet $\oplus$ singlet)
diagonalizes $\Omega_{12}$:
\[
 \Omega_{12}\big|_{V_1}=\tfrac{1}{2}\,\id,
 \qquad
 \Omega_{12}\big|_{V_0}=-\tfrac{3}{2}\,\id.
\]
Indeed, $\frac{1}{2}(j_{\mathrm{tot}}(j_{\mathrm{tot}}+1)-2\cdot\frac{1}{2}\cdot\frac{3}{2})$
gives $\frac{1}{2}$ for $j_{\mathrm{tot}}=1$ and $-\frac{3}{2}$ for $j_{\mathrm{tot}}=0$.
\end{computation}

\begin{computation}[$\Conf_2(\CC)$ monodromy at general level;
\ClaimStatusProvedHere]
\label{comp:kz-monodromy-conf2}
\index{KZ connection!monodromy}
The KZ connection at level~$k\neq -2$ on
$\Conf_2(\CC)$ with two fundamental insertions is
\[
 \nabla \;=\; d \;-\; \frac{1}{k+2}\,\Omega_{12}\,d\log(z_1-z_2).
\]
Since $\Omega_{12}$ is diagonalizable, parallel transport around
the generator $\gamma$ of $\pi_1(\Conf_2(\CC))\cong\ZZ$
(the loop in which $z_1$ encircles $z_2$ counterclockwise) gives
\begin{equation}\label{eq:conf2-monodromy}
 M_\gamma
 \;=\;
 \exp\!\bigl(-2\pi i\cdot\tfrac{1}{k+2}\,\Omega_{12}\bigr)
 \;=\;
 \begin{cases}
 e^{-\pi i/(k+2)}\,\id & \text{on } V_1, \\[3pt]
 e^{3\pi i/(k+2)}\,\id & \text{on } V_0.
 \end{cases}
\end{equation}
Setting $q=e^{i\pi/(k+2)}$, the eigenvalues are $q^{-1}$ (triplet) and $q^3=-q^{-1}\cdot q^{k+2}=-q^{-1}$
when $q^{k+2}=1$... more precisely, $e^{3\pi i/(k+2)}=q^3$.
These are exactly the eigenvalues of the $R$-matrix of $U_q(\mathfrak{sl}_2)$
on $V\otimes V$:
\[
 \check{R}\big|_{V_1}=q,\qquad \check{R}\big|_{V_0}=-q^{-1},
\]
matching $M_\gamma = q^{-\Omega_{12}} = P\circ \check{R}^{-1}$
(with $P$ the transposition) up to the standard framing scalar.
At $k=1$ one recovers $q=e^{i\pi/3}$, $M_\gamma=\operatorname{diag}(-1,\,q^{-1})$
on the singlet/triplet decomposition, confirming the
$\widehat{\mathfrak{sl}}_2$ level-$1$ WZW fusion rules.
\end{computation}

\begin{computation}[Formal $R$-matrix expansion; \ClaimStatusProvedHere]
\label{comp:formal-R-expansion}
\index{R-matrix@$R$-matrix!formal expansion}
The formal monodromy operator, equivalently the asymptotic $R$-matrix,
admits the $z^{-1}$-expansion (with $\Omega$ denoting the Sugawara-normalized
residue tensor $k\,\Omega_{\mathrm{Cas}}/(k+h^\vee)$ absorbing the level prefix and the level-explicit form at~\eqref{eq:R-level-explicit} below)
\begin{equation}\label{eq:R-formal-expansion}
 R(z) \;=\; 1 + \frac{\Omega}{z} + \frac{\Omega^2}{2z^2} + \frac{\Omega^3}{6z^3}+\cdots
 \;=\; \exp\!\Bigl(\frac{\Omega}{z}\Bigr)\quad (\text{leading order}).
\end{equation}
More precisely, at level~$k$ the ratio $\Omega/(k+2)$ replaces $\Omega$
and corrections of order $O(z^{-3})$ arise from the non-abelian exponentiation
(Magnus series). On $V\otimes V$ one can diagonalize:
\[
 R(z)\big|_{V_1} = e^{1/(2z)}\bigl(1+O(z^{-3})\bigr),
 \qquad
 R(z)\big|_{V_0} = e^{-3/(2z)}\bigl(1+O(z^{-3})\bigr).
\]
The leading-order formula
\begin{equation}\label{eq:R-level-explicit}
 R(z)=1+k\,\Omega/z+O(z^{-2})
\end{equation}
at level~$k$ identifies $k\,\Omega$ as the collision residue, consistent with
Lemma~\ref{lem:inf-braid}.
\end{computation}

\begin{computation}[$\Conf_3^{\mathrm{ord}}(\CC)$ monodromy and Yang--Baxter;
\ClaimStatusProvedHere]
\label{comp:conf3-yang-baxter}
\index{KZ connection!three-point!Yang--Baxter}
On $\Conf_3^{\mathrm{ord}}(\CC)$ with three fundamental insertions, the KZ connection is
\[
 \nabla_3 \;=\; d - \frac{1}{k+2}\Bigl(
 \Omega_{12}\,d\log z_{12}+\Omega_{13}\,d\log z_{13}+\Omega_{23}\,d\log z_{23}
 \Bigr),
\]
where $z_{ij}:=z_i-z_j$. The fundamental group
$\pi_1(\Conf_3^{\mathrm{ord}}(\CC))$ is the pure braid group $P_3$,
generated by the three elementary braids $\sigma_{ij}^2$
(particle~$i$ encircling particle~$j$). The monodromy representation
\[
 \rho\colon P_3\;\longrightarrow\; \mathrm{GL}(V^{\otimes 3})
\]
assigns to each generator
$\rho(\sigma_{ij}^2) = \exp\!\bigl(-2\pi i\cdot\tfrac{1}{k+2}\,\Omega_{ij}\bigr)$,
acting in the $(i,j)$-tensor factors.

The pure braid relation $(\sigma_{12}\sigma_{23})^3=(\sigma_{12}\sigma_{23}\sigma_{12})^2$
is verified at the infinitesimal level by the identity
\begin{equation}\label{eq:inf-braid-yb}
 [\Omega_{12},\Omega_{13}+\Omega_{23}]=0,
\end{equation}
which is the infinitesimal braid relation (Lemma~\ref{lem:inf-braid}).
Exponentiating, the monodromies satisfy the quantum Yang--Baxter equation
on $V^{\otimes 3}$:
\begin{equation}\label{eq:qybe-conf3}
 R_{12}R_{13}R_{23}=R_{23}R_{13}R_{12},
\end{equation}
where $R_{ij}:=\exp\!\bigl(-2\pi i\cdot\frac{1}{k+2}\,\Omega_{ij}\bigr)$.
This is a direct consequence of the flatness theorem
(Theorem~\ref{thm:strict-flatness}): parallel transport around two
homotopic loops in $\Conf_3(\CC)$ must agree, and the two decompositions
of the full braid into elementary braids yield the two sides
of~\eqref{eq:qybe-conf3}.
\end{computation}

\begin{computation}[Logarithmic extension to $\FM_2(\PP^1\mid\infty)$;
\ClaimStatusProvedHere]
\label{comp:log-extension-fm2}
\index{Fulton--MacPherson!logarithmic extension}
The space $\FM_2(\PP^1\mid\infty)$ compactifies $\Conf_2(\PP^1\setminus\{\infty\})$
by adjoining a boundary divisor $D_{\mathrm{coll}}$ (the collision locus $z_1=z_2$)
and a boundary divisor $D_\infty$ (a point escaping to $\infty$).
In center-of-mass / relative coordinates $w=(z_1+z_2)/2$, $u=z_1-z_2$:
\begin{enumerate}[label=\textup{(\alph*)}]
\item \emph{Collision boundary} ($u\to 0$).
 The connection becomes
 $\nabla = d - \frac{1}{k+2}\,\Omega_{12}\,\frac{du}{u} + (\text{regular in } u)$.
 By Proposition~\ref{prop:strict-log-ext}, the residue along $D_{\mathrm{coll}}$
 is $\Res_{D_{\mathrm{coll}}}(\nabla)=\frac{1}{k+2}\,\Omega_{12}$.
 The local monodromy around $D_{\mathrm{coll}}$ is
 $\exp\!\bigl(-\frac{2\pi i}{k+2}\,\Omega_{12}\bigr)$ as computed in~\eqref{eq:conf2-monodromy}.
\item \emph{Infinity boundary} ($w\to\infty$ with $u$ bounded).
 In the coordinate $s=1/w$, the connection is regular at $s=0$
 because the KZ form $d\log(z_1-z_2)=d\log u$ has no pole in~$w$.
 Hence the local monodromy around $D_\infty$ is trivial.
\item \emph{Corner stratum} ($u\to 0$ and $w\to\infty$ simultaneously).
 Both $z_1,z_2\to\infty$ while colliding.
 In the chart $(s,t)=(1/w,\,u\cdot w)$, the collision scale coordinate
 is $t$ and the connection form is
 $\frac{1}{k+2}\,\Omega_{12}\,\frac{dt}{t}+O(1)$,
 confirming logarithmic regularity at the corner.
\end{enumerate}
Thus $\nabla$ extends to a logarithmic connection on $\FM_2(\PP^1\mid\infty)$
with a single nontrivial residue $\frac{1}{k+2}\Omega_{12}$
along the collision divisor.
\end{computation}

\begin{computation}[Genus-$1$ KZB connection and modular monodromy;
\ClaimStatusProvedHere]
\label{comp:genus1-kzb}
\index{KZB connection!genus one!sl2@$\mathfrak{sl}_2$}
On an elliptic curve $E_\tau=\CC/(\ZZ+\ZZ\tau)$ ($\Im\tau>0$) with two
marked points at positions $0$ and $z$,
the genus-$0$ propagator $1/(z_1-z_2)$ is replaced by the
Weierstrass $\zeta$-function (or equivalently $\partial_z\log\theta_1$).
The KZB connection on $E_\tau\setminus\{0\}$ with fiber $V\otimes V$ is
\begin{equation}\label{eq:kzb-connection}
 \nabla_\tau \;=\; d_z \;-\; \frac{1}{k+2}\,\Omega_{12}\,
 \partial_z\!\log\theta_1(z;\tau)\,dz,
\end{equation}
where $\theta_1(z;\tau)=2\sum_{n=0}^{\infty}(-1)^n q^{(n+1/2)^2}\sin((2n+1)\pi z)$,
$q=e^{i\pi\tau}$, is the odd Jacobi theta function vanishing at $z=0$.

\medskip\noindent\textbf{Near $z=0$.}\;
Since $\theta_1(z;\tau)=2\pi z\,\eta(\tau)^3(1+O(z^2))$,
one has $\partial_z\log\theta_1=1/z+O(z)$, and the connection reduces to
\[
 \nabla_\tau = d_z - \frac{1}{k+2}\,\Omega_{12}\,\frac{dz}{z}+O(1),
\]
recovering the genus-$0$ KZ connection locally. The local monodromy around
$z=0$ is therefore the same as~\eqref{eq:conf2-monodromy}.

\medskip\noindent\textbf{$A$-cycle monodromy} ($z\mapsto z+1$).\;
The quasi-periodicity $\theta_1(z+1;\tau)=-\theta_1(z;\tau)$ gives
$\partial_z\log\theta_1(z+1;\tau)=\partial_z\log\theta_1(z;\tau)$.
Hence the connection form is strictly periodic under $z\mapsto z+1$
and the $A$-cycle monodromy is the identity:
\begin{equation}\label{eq:A-cycle-monodromy}
 M_A = \id_{V\otimes V}.
\end{equation}

\medskip\noindent\textbf{$B$-cycle monodromy} ($z\mapsto z+\tau$).\;
The quasi-periodicity
$\theta_1(z+\tau;\tau)=-e^{-i\pi\tau-2\pi iz}\theta_1(z;\tau)$ yields
\[
 \partial_z\log\theta_1(z+\tau;\tau)
 = \partial_z\log\theta_1(z;\tau)-2\pi i.
\]
Under transport $z\mapsto z+\tau$, the connection form shifts:
\[
 \nabla_\tau\big|_{z+\tau}
 =
 d_z - \frac{1}{k+2}\,\Omega_{12}\bigl(\partial_z\log\theta_1(z;\tau)-2\pi i\bigr)dz
 =
 \nabla_\tau + \frac{2\pi i}{k+2}\,\Omega_{12}\,dz.
\]
The $B$-cycle monodromy therefore acquires an abelian twist
controlled by the integrated correction $\int_0^\tau \frac{2\pi i}{k+2}\,\Omega_{12}\,dz
= \frac{2\pi i\tau}{k+2}\,\Omega_{12}$:
\begin{equation}\label{eq:B-cycle-monodromy}
 M_B = \exp\!\Bigl(\frac{2\pi i\tau}{k+2}\,\Omega_{12}\Bigr)
 \cdot M_{\mathrm{loc}},
\end{equation}
where $M_{\mathrm{loc}}$ is the genus-$0$ local monodromy from~\eqref{eq:conf2-monodromy}
(arising from the pole at $z=0$ enclosed by the $B$-cycle).
On the Clebsch--Gordan components:
\[
 M_B\big|_{V_1} = e^{\pi i\tau/(k+2)}\cdot e^{-\pi i/(k+2)}
 = e^{\pi i(\tau-1)/(k+2)},
 \qquad
 M_B\big|_{V_0} = e^{-3\pi i\tau/(k+2)}\cdot e^{3\pi i/(k+2)}
 = e^{-3\pi i(\tau-1)/(k+2)}.
\]

\medskip\noindent\textbf{Modular $R$-matrix.}\;
Define the genus-$1$ $R$-matrix as the ratio of the connection
form~\eqref{eq:kzb-connection} to its genus-$0$ counterpart:
\begin{equation}\label{eq:genus1-r-matrix}
 r_1(z;\tau) \;:=\; \frac{k}{k+2}\,\Omega\cdot \partial_z\!\log\theta_1(z;\tau).
\end{equation}
Expanding in $z$:
\[
 r_1(z;\tau)
 = \frac{k\,\Omega}{(k+2)z}
 - \frac{k\,\Omega}{k+2}\sum_{n=1}^{\infty}2\,
 \frac{q^{2n}}{1-q^{2n}}\cdot 2\pi i\cos(2\pi nz)
 = \frac{k\,\Omega}{(k+2)z}
 + \frac{k\,\Omega}{k+2}\bigl(2G_2(\tau)-\wp(z;\tau)-\pi^2/3\bigr),
\]
where the Eisenstein correction $G_2(\tau)$ accounts for the
(conditionally convergent) sum and $\wp$ is the Weierstrass
$\wp$-function. (The genus-$0$ collision residue is $r(z) = k\,\Omega/((k+h^\vee)z)$; the intrinsic
KZB connection~\eqref{eq:kzb-connection} retains the standard
Sugawara normalization $(k+h^\vee)^{-1}$, so the identification of
$r_1$ with the connection form carries the rescaling factor
$k/(k+h^\vee)$. At $k=0$ both the connection trivializes in the
collision sector and $r_1$ vanishes, consistent with
Theorem~\ref{thm:affine-monodromy-identification}.) Thus the
genus-$1$ $R$-matrix is the genus-$0$
rational $R$-matrix $k\,\Omega/((k+2)z)$ dressed by the elliptic
Weierstrass function, exactly as predicted by the shadow obstruction tower:
the genus-$1$ period correction to the Casimir is controlled by
$G_2(\tau)$ and reproduces the curvature $\kappa\cdot\omega_1$
of the modular shadow at genus~$1$
(cf.\ Theorem~\ref{thm:synthesis}\textup{(i)}).
\end{computation}

\begin{remark}[Consistency with the modular shadow obstruction tower]
\label{rem:kzb-shadow-consistency}
The genus-$1$ connection~\eqref{eq:kzb-connection} is the $n=2$ projection
of the shadow connection $\nabla_{g,n}=d-\mathrm{Sh}_{g,n}(\Theta_A)$
from the synthesis theorem, specialized to $\mathfrak{g}=\mathfrak{sl}_2$
and $g=1$. The Eisenstein correction $G_2(\tau)$ in $r_1$
is the genus-$1$ curvature $\kappa(\widehat{\mathfrak{sl}}_2)
=\dim\mathfrak{g}\cdot(k+h^\vee)/(2h^\vee)=3(k+2)/4$
times the period form $\omega_1$, in agreement with the formula
$\kappa(\cA)\cdot\omega_g$ established in Volume~I.
\end{remark}

\section{Monodromy equals spectral $R$-matrix: the affine lineage}\label{sec:affine-monodromy-identification}

The computations above identify the reduced HT monodromy for $\widehat{\fg}_k$ with the eigenvalues of the quantum group $R$-matrix. The proof of this identification as a theorem, valid for every simple Lie algebra $\fg$ at generic level, by combining three inputs: one-loop exactness of the affine half-space BV theory, the Drinfeld--Kohno theorem, and the spectral $R$-matrix of Section~\ref{sec:spectral_braiding}.

The key mechanism is that one-loop exactness \emph{collapses the $A_\infty$ tower}: the higher collision homotopies $h_S$ ($|S| \geq 3$) of Section~\ref{sec:ainfty} vanish after reduction to degree-zero states, because the tree-level BRST complex is a regular-sequence Koszul model (Proposition~\ref{prop:regular-sequence}). The reduced logarithmic connection is therefore determined entirely by its binary residues, and these are the Casimir tensors of the KZ connection.

\subsection{One-loop collapse and the KZ identification}

Let $\fg$ be a simple Lie algebra with dual Coxeter number $h^\vee$, Killing form $\kappa_\fg$, and split Casimir $\Omega \in \fg \otimes \fg$. Let $k \in \CC$ with $k \neq -h^\vee$. Consider the affine Kac--Moody vertex algebra $V^k(\fg)$ as a logarithmic $\SCchtop$-algebra via the physics bridge (Theorem~\ref{thm:physics-bridge}).

\begin{proposition}[Affine line operators are resolved in degree $0$; \ClaimStatusProvedHere]
\label{prop:affine-resolved}
Let $V_1, \dots, V_n$ be finite-dimensional $\fg$-modules, viewed as evaluation modules for the dg-shifted Yangian $\cA^!_{\mathrm{line}}$ at spectral positions $z_1, \dots, z_n$ (an \emph{evaluation module} is a finite-dimensional $\fg$-module $V$ pulled back to a $Y_\hbar(\fg)$-module via the evaluation homomorphism $\mathrm{ev}_z: Y_\hbar(\fg) \to U(\fg)$ sending $T(u) \mapsto 1 + \hbar J/(u-z)$). Then the tree-level BRST complex $(E_n, Q_{0,n})$ is resolved in degree $0$: $H^j(E_n, Q_{0,n}) = 0$ for $j \neq 0$.
\end{proposition}

\begin{proof}
The affine half-space BV quantization is one-loop exact (Theorem~\ref{thm:affine-half-space-bv}). By the one-loop Koszul criterion (Theorem~\ref{thm:one-loop-koszul}), the pre-reduction boundary algebra $V^k(\fg)$ is chirally Koszul: the bar cohomology is concentrated and the transferred $A_\infty$ operations satisfy $m_k^{\mathrm{aff}} = 0$ for $k \geq 3$. Consequently the tree-level line complex for evaluation modules is a Koszul complex of the type in Proposition~\ref{prop:regular-sequence}, with the regular sequence given by the spectral evaluation constraints. The concentration in degree $0$ follows.
\end{proof}

\begin{theorem}[Reduced HT connection equals KZ connection; \ClaimStatusProvedHere]
\label{thm:reduced-equals-kz}
In the setting of Proposition~\textup{\ref{prop:affine-resolved}}, the reduced logarithmic flat connection on $H_n = H^0(E_n[[\h]], Q_n)$ is the Knizhnik--Zamolodchikov connection:
\[
 \A_n \;=\; \frac{1}{k + h^\vee}\sum_{1 \leq i < j \leq n} \rho_{ij}(\Omega)\, d\log(z_i - z_j),
\]
where $\rho_{ij}$ acts in the $i$-th and $j$-th tensor factors of\/ $V_1 \otimes \cdots \otimes V_n$.
\end{theorem}

\begin{proof}
By Proposition~\ref{prop:affine-resolved}, the tree-level line complex is resolved in degree~$0$. The unconditionality theorem (Theorem~\ref{thm:unconditionality}) therefore applies: the canonical $\h$-adic retract exists, and the reduced superconnection is an ordinary logarithmic flat connection
\[
 \A_n = p_n A_n^{(1)} i_n
\]
by one-form rigidity (Theorem~\ref{thm:one-form-rigidity}).

It remains to identify the binary residues. By Corollary~\ref{cor:reduced-residues}, the residue of $\A_n$ along the collision divisor $D_{\{i,j\}}$ is
\[
 \Res_{D_{\{i,j\}}}(\A_n) = p_n \Res_{D_{\{i,j\}}}(A_n^{(1)}) \, i_n.
\]
The raw $1$-form residue $\Res_{D_{\{i,j\}}}(A_n^{(1)})$ is the binary collision kernel $r_{ij}(z_{ij})$ evaluated at $z_{ij} = 0$, which by the Laplace formula (Proposition~\ref{prop:field-theory-r}) is the leading pole of
\[
 r(z) = \int_0^\infty e^{-\lambda z} \{\cdot{}_\lambda \cdot\}\, d\lambda.
\]
For the affine Kac--Moody algebra $V^k(\fg)$, the $\lambda$-bracket on generators is
\[
 \{J^a{}_\lambda J^b\} = f^{ab}_c J^c + k \kappa_\fg^{ab} \lambda.
\]
The Laplace transform gives the full collision kernel
\[
 r(z) = \frac{1}{z}\Big(\sum_{a,b} f^{ab}_c J^c \otimes E_{ab}\Big) + \frac{k\,\kappa_\fg}{z^2} + O(z^{-3}).
\]
\noindent
(This is the Laplace kernel $r^L(z)$; the bar-theoretic collision
residue $r^{\mathrm{coll}}(z)$, obtained after $d\log$ absorption,
has pole orders one lower.)
On evaluation modules, where $J^a$ acts as the finite-dimensional representation matrix $\rho(X^a)$, the residue at $z = 0$ is determined by the simple-pole coefficient. The structure-constant tensor $\sum_{a,b} f^{ab}_c \rho(X^c) \otimes E_{ab}$ equals the adjoint action of the split Casimir~$\Omega = \sum_a X^a \otimes X_a$ (with indices raised by the finite Killing form~$\kappa_\fg$):
\[
 \Res_{z=0}\, r(z)\,dz \;=\; \Omega_{ij}.
\]
The factor $(k + h^\vee)^{-1}$ arises from the passage to the \emph{reduced} connection via the projection~$p_n$. The bar-complex differential normalises binary operations by the affine bilinear form $(k + h^\vee)\kappa_\fg$ (the $h^\vee$-shift being the normal-ordering correction in the Segal--Sugawara vector $T = \tfrac{1}{2(k+h^\vee)} \sum{:}J^a J_a{:}$, which ensures $T(z)J^a(w) \sim J^a(w)/(z-w)^2$). The projection $p_n$ to degree-zero states absorbs this normalisation, giving the KZ residue $\Omega_{ij}/(k + h^\vee)$.

Since the affine algebra is Koszul ($m_k = 0$ for $k \geq 3$), there are no higher collision corrections to the residue: the $A_\infty$ Yang--Baxter defect (Theorem~\ref{thm:first-obstruction}) vanishes identically, and the higher support homotopies $h_S$ ($|S| \geq 3$) can be taken to be zero. After projection to degree-zero states, the regular terms in $r(z)$ beyond the simple pole act on higher-degree components that are killed by $p_n$. Thus
\[
 \A_n = \frac{1}{k + h^\vee}\sum_{i < j} \rho_{ij}(\Omega)\, d\log(z_i - z_j),
\]
which is precisely the KZ connection at level $k$.
\end{proof}

\subsection{The identification theorem}

\begin{theorem}[Monodromy equals spectral $R$-matrix for the affine lineage; \ClaimStatusProvedHere]
\label{thm:affine-monodromy-identification}
\index{monodromy!affine identification|textbf}
\index{R-matrix@$R$-matrix!monodromy identification|textbf}
Let $\fg$ be a simple Lie algebra and $k \in \CC \setminus \{-h^\vee\}$ generic. Let $V_1, \dots, V_n$ be finite-dimensional $\fg$-modules. Then:
\begin{enumerate}[label=\textup{(\roman*)}]
 \item The monodromy of the reduced HT logarithmic connection on $\Conf_n(\CC)$ equals the monodromy of the KZ connection at level $k$:
 \[
 \rho_n^{\mathrm{HT}} = \rho_n^{\mathrm{KZ}} \colon \pi_1(\Conf_n(\CC)) \longrightarrow \Aut(V_1 \otimes \cdots \otimes V_n).
 \]
 \item By the Drinfeld--Kohno theorem, this monodromy is compared with the corresponding quantum-group braid representation at $q = e^{i\pi/(k + h^\vee)}$: the representation $\rho_n^{\mathrm{HT}}$ is equivalent to the braided tensor product representation of\/ $\mathrm{Rep}_q(\fg)$ on $V_1 \otimes \cdots \otimes V_n$.
 \item The spectral $R$-matrix $R_{V_i, V_j}(z)$ of Section~\textup{\ref{sec:spectral_braiding}}, restricted to evaluation modules, equals the quantum group $R$-matrix: on the reduced evaluation comparison surface, the braided monoidal structure of the reduced HT line category agrees with the braided monoidal structure of $\mathrm{Rep}_q(\fg)$.
\end{enumerate}
\end{theorem}

\begin{proof}
\textbf{(i)} is immediate from Theorem~\ref{thm:reduced-equals-kz}: the reduced HT connection \emph{is} the KZ connection, so their monodromies are identical.

\medskip
\noindent\textbf{(ii).} The Drinfeld--Kohno theorem (\cite{Dri89}, Theorem~3.2; \cite{KL93}) compares the monodromy representation of the KZ connection at level $k$, acting on tensor products of finite-dimensional $\fg$-modules, with the braided tensor product representation of $\mathrm{Rep}_q(\fg)$ at $q = e^{i\pi/(k + h^\vee)}$. The braiding $\sigma_{ij}^2 \mapsto \exp(-2\pi i \cdot \Omega_{ij}/(k + h^\vee))$ on elementary generators matches the universal $R$-matrix of $U_q(\fg)$ acting on $V_i \otimes V_j$. Since $\rho_n^{\mathrm{HT}} = \rho_n^{\mathrm{KZ}}$, the same comparison holds for the HT monodromy.

\medskip
\noindent\textbf{(iii).} We must show that the spectral $R$-matrix
$R_{V_i,V_j}(z)$ of Definition~\ref{def:spectral-braiding},
restricted to evaluation modules, equals the quantum group
$R$-matrix. The argument has three sub-steps: reduction of the
$A_\infty$ tower, identification of the classical $r$-matrix, and
exponentiation.

\smallskip\noindent
\emph{Step (iii-a): One-loop collapse on evaluation modules.}
The spectral $R$-matrix is defined by configuration integrals
over $\FM_2(\CC)$ (Theorem~\ref{thm:braided-category}),
encoding the braiding on the boundary line category
$\cC_\partial$. In full generality, the braiding receives
contributions from all $A_\infty$ operations $m_k$ with
$k \geq 2$, corresponding to multi-point configuration
integrals on $\FM_k(\CC)$ via the iterated integral formula
for the path-ordered exponential.

For the affine algebra $V^k(\fg)$ at generic level $k$,
Theorem~\ref{thm:affine-half-space-bv} proves one-loop
exactness: the operations $m_k = 0$ for $k \geq 3$
(only trees and one-loop graphs contribute, and for the
affine algebra the one-loop graph determines $m_2$ exactly
while $m_k = 0$ for $k \geq 3$ by the graph-counting argument
in the proof of Theorem~\ref{thm:affine-half-space-bv}). Therefore the
path-ordered exponential reduces to the standard ordered
exponential of the flat connection defined by $m_2$ alone:
\[
 R(z) = \mathrm{P}\!\exp\!\int_\gamma r(\zeta)\, d\zeta,
\]
where $\gamma$ is the half-monodromy path in $\Conf_2(\CC)$
and $r(z)$ is the infinitesimal $r$-matrix extracted from $m_2$.
No higher $m_k$ corrections appear.

\smallskip\noindent
\emph{Step (iii-b): Classical $r$-matrix equals Casimir kernel.}
By Proposition~\ref{prop:field-theory-r}, the infinitesimal
spectral $r$-matrix is
\[
 r(z) = \int_0^\infty e^{-\lambda z}
 \{\cdot{}_\lambda\cdot\}\, d\lambda.
\]
For $V^k(\fg)$, the $\lambda$-bracket is
$\{J^a {}_\lambda J^b\} = f^{ab}_c J^c + k\,\kappa_\fg^{ab}\lambda$
(where $\kappa_\fg$ is the Killing form and $f^{ab}_c$ are the
structure constants). Evaluating the Laplace transform:
\[
 r(z) = \frac{\sum_{a,b} J^a \otimes J^b \cdot (\kappa_\fg)_{ab}}{z}
 + \frac{k\,\kappa_\fg}{z^2}
 = \frac{\Omega}{z} + \frac{k\,\kappa_\fg}{z^2},
\]
where $\Omega = \sum_a J^a \otimes J_a \in \fg \otimes \fg$
is the quadratic Casimir tensor. On evaluation modules
$V_i, V_j$, the representation $\rho_{ij}(\Omega)$ acts as
a finite-rank operator, and the $z^{-2}$ term is a scalar
(proportional to the identity on $V_i \otimes V_j$). The
scalar term does not contribute to the braiding (it exponentiates
to a scalar phase), so the braiding is governed by
$r(z) = k\,\Omega/z$ at level~$k$ up to scalars.

\smallskip\noindent
\emph{Step (iii-c): Exponentiation to KZ monodromy.}
The flat connection with connection form
$\Omega_{ij}\, d\log(z_i - z_j)/(k + h^\vee)$ is precisely
the KZ connection of Theorem~\ref{thm:reduced-equals-kz}.
Its half-monodromy (the braiding operator obtained by
analytically continuing $z_i$ around $z_j$ along a half-circle)
is
\[
 R_{ij} = \exp\!\Bigl(\frac{-\pi i\, \rho_{ij}(\Omega)}
 {k + h^\vee}\Bigr).
\]
This follows from the standard computation: the
path-ordered exponential of $r(\zeta)\,d\zeta$ along the
half-circle $\zeta = z_{ij} e^{i\theta}$,
$\theta \in [0,\pi]$, gives
$\mathrm{P}\!\exp(\int_0^\pi (k\,\Omega/z_{ij} e^{i\theta})
\cdot iz_{ij} e^{i\theta}\, d\theta)
= \exp(i\pi\,k\, \Omega)$, with the normalization
$1/(k+h^\vee)$ from Theorem~\ref{thm:reduced-equals-kz}.

By the Drinfeld--Kohno theorem (part~(ii)), this half-monodromy
is exactly the image of the universal $R$-matrix of $U_q(\fg)$
on $V_i \otimes V_j$. Therefore
\[
 R_{V_i,V_j}^{\mathrm{spectral}}
 = R_{V_i,V_j}^{\mathrm{KZ}}
 = R_{V_i,V_j}^{U_q(\fg)}.
\]

\smallskip\noindent
\textbf{Combining (i)--(iii):} On evaluation modules, the
spectral $R$-matrix equals the KZ monodromy (by the chain
(iii-a)$\to$(iii-b)$\to$(iii-c)), which equals the quantum group
$R$-matrix (by (ii)). The full Yang--Baxter structure is preserved:
the spectral $R$-matrix satisfies YBE by
Theorem~\ref{thm:YBE}, and the quantum group $R$-matrix satisfies
YBE by the quasi-triangular structure of $U_q(\fg)$; the
identification of the two on evaluation modules is compatible
because the braided monoidal functor
$\mathrm{ev} \colon \mathrm{Rep}_q(\fg) \to \cC_\partial^{\mathrm{red}}$
preserves the braiding (this is the content of the identification).
The associator $\Phi_{\mathrm{HT}}$ of
Definition~\ref{def:ht-assoc} likewise reduces to the Drinfeld
associator $\Phi_{\mathrm{KZ}}$ by the same one-loop collapse:
the three-point $\FM_3$-integral reduces to the KZ three-point
connection, whose monodromy is $\Phi_{\mathrm{KZ}}$. The pentagon
and hexagon identities
(Theorems~\ref{thm:pentagon} and~\ref{thm:hexagon}) are therefore
the pentagon and hexagon of $\mathrm{Rep}_q(\fg)$, and on the
evaluation-module sector the two braided monoidal categories
coincide.
\end{proof}

\begin{remark}[Unconditional status]
\label{rem:affine-unconditional}
Every input to Theorem~\ref{thm:affine-monodromy-identification} is unconditional. The one-loop exactness of the affine half-space BV theory is proved in Theorem~\ref{thm:affine-half-space-bv}. The Koszulness of $V^k(\fg)$ follows from Theorem~\ref{thm:one-loop-koszul}. The unconditionality theorem (Theorem~\ref{thm:unconditionality}) guarantees the reduced monodromy machine. The Drinfeld--Kohno theorem is a classical result. No standing hypothesis remains: the result holds for any simple $\fg$ and any generic level $k$.
\end{remark}

\begin{remark}[Spectral-to-categorical identification]
\label{rem:spectral-categorical-crossref}
Theorem~\ref{thm:affine-monodromy-identification} is the affine instance of the general spectral-to-categorical bridge (Proposition~\ref{prop:spectral-to-categorical}): the spectral $R$-matrix $R(z)$, integrated along a half-circle exchange path in $\Conf_2(\CC)$, yields the categorical braiding of $\mathrm{Rep}_q(\fg)$.  This is the content of the Drinfeld--Kohno theorem.  See Remark~\ref{rem:spectral-vs-categorical} for the general relationship between spectral and categorical braidings.
\end{remark}

\subsection{The $\mathfrak{sl}_3$ verification: cubic Casimir and the associator}
\label{subsec:sl3-verification}

\providecommand{\fsl}{\mathfrak{sl}}

Theorem~\ref{thm:affine-monodromy-identification} holds for all simple $\fg$. For $\fg = \fsl_2$, the monodromy is controlled entirely by the quadratic Casimir $C_2$. The first genuinely higher-rank phenomenon appears for $\fg = \fsl_3$, where the cubic Casimir $C_3$ (from the totally symmetric invariant tensor $d_{abc}$) enters the monodromic structure. The following proposition records the explicit verification.

\begin{proposition}[$\fsl_3$ monodromy identification on evaluation modules; \ClaimStatusProvedHere]
\label{prop:sl3-monodromy-verification}
\index{sl3 monodromy@$\mathfrak{sl}_3$ monodromy}
Let $\fg = \fsl_3$ with dual Coxeter number $h^\vee = 3$, and let $V = V_{(1,0)} = \CC^3$ be the fundamental representation. Set $q = e^{i\pi/(k+3)}$.
\begin{enumerate}[label=\textup{(\roman*)}]
 \item \emph{Braiding on $V \ot V$}: The Clebsch--Gordan decomposition $V \ot V \cong V_{(2,0)} \oplus V_{(0,1)}$ gives split Casimir eigenvalues $\Omega|_{V_{(2,0)}} = 2/3$ and $\Omega|_{V_{(0,1)}} = -4/3$. The monodromy of the bar-complex Steinberg connection $\nabla_{\fS} = d + \Omega\,dz/((k+3)z)$ matches the $U_q(\fsl_3)$ $\check{R}$-matrix:
 \[
 \mathrm{Mon}\big|_{V_{(2,0)}} = q^{4/3},
 \qquad
 \mathrm{Mon}\big|_{V_{(0,1)}} = q^{-8/3}.
 \]
 The cubic Casimir does not appear: $C_2$ alone distinguishes the two summands.

 \item \emph{Associator on $V^{\ot 3}$}: The tensor cube decomposes as $V^{\ot 3} \cong V_{(3,0)} \oplus 2\cdot V_{(1,1)} \oplus V_{(0,0)}$. The two copies of the adjoint $V_{(1,1)}$ arise from different intermediate channels \textup{(}via $V_{(2,0)}$ and $V_{(0,1)}$ respectively in the $(12)$-grouping\textup{)}. At order $\hbar^3 = 1/(k+3)^3$ in the Drinfeld--KZ associator, the antisymmetric double commutator $\mathcal{A} = [\Omega_{12}, [\Omega_{12}, \Omega_{23}]] - [\Omega_{23}, [\Omega_{23}, \Omega_{12}]]$ receives a contribution from the $d$-tensor of $\fsl_3$ via the Fierz identity
 \[
 f_{ace}\,f_{bde} = \tfrac{2}{3}(\delta_{ab}\delta_{cd} - \delta_{ad}\delta_{bc}) + d_{ace}\,d_{bde} - d_{ade}\,d_{bce}.
 \]
 The $d$-tensor piece $\mathcal{D}_{123}^{(-)}$ produces off-diagonal mixing between the two copies of $V_{(1,1)}$, proportional to the cubic Casimir difference $\Delta c_3 = c_3(2,0) - c_3(0,1) = 5$. For $\fsl_2$, $d_{abc} = 0$ and no such mixing occurs.

 \item \emph{Bar-complex identification}: By one-loop exactness of $V^k(\fsl_3)$ at generic $k$ \textup{(}Theorem~\textup{\ref{thm:one-loop-koszul})}, the $A_\infty$ operations $m_k = 0$ for $k \geq 3$, so the reduced HT connection is exactly the KZ connection. Therefore the bar-complex associator reproduces the Drinfeld--KZ associator identically, including the cubic Casimir contribution. By the Drinfeld--Kohno theorem, the monodromy identifies with the $U_q(\fsl_3)$ braiding and associator on evaluation modules.
\end{enumerate}
\end{proposition}

\begin{proof}
The computation is a direct verification. The split Casimir $\Omega = \sum_{i \neq j} E_{ij} \ot E_{ji} + \sum_{a,b} (A^{-1})_{ab} H_a \ot H_b$ has the stated eigenvalues by the formula $\Omega|_{V_\nu} = \frac{1}{2}(c_2(\nu) - c_2(\lambda) - c_2(\mu))$. For part~(ii), the Fierz decomposition is standard; the key point is that $\mathcal{D}_{123}^{(-)} \neq 0$ because the totally symmetric tensor $d_{abc}$ of $\fsl_3$ is nonvanishing. The off-diagonal matrix element on the multiplicity space of $V_{(1,1)}$ is proportional to $c_3(2,0) - c_3(0,1) = 35/9 + 10/9 = 5$. For part~(iii), one-loop exactness gives $m_k = 0$ for $k \geq 3$, and the reduced connection $\nabla_{\fS} = d + \Omega\,dz/((k+3)z)$ is the KZ connection for $\fsl_3$ at level $k$. The identification with $U_q(\fsl_3)$ then follows from Theorem~\ref{thm:affine-monodromy-identification} applied to $\fg = \fsl_3$.
\end{proof}

\begin{remark}[The cubic Casimir enters through the associator, not the $R$-matrix]
\label{rem:cubic-through-associator}
On $V \ot V$, the $\fsl_3$ braiding looks identical to the $\fsl_2$ braiding: both are diagonal on Clebsch--Gordan summands with eigenvalues determined by $C_2$. The genuinely new higher-rank phenomenon is visible only in the associator on $V^{\ot 3}$ (or in the braiding on tensor products with outer multiplicities). The cubic Casimir controls the off-diagonal mixing between multiplicity copies; the quadratic Casimir controls the diagonal part. This is the first explicit verification that the bar-complex monodromy machine correctly captures higher Casimir contributions.
\end{remark}

\subsection{Corollaries}

\begin{corollary}[Reduced HT line category for the affine lineage; \ClaimStatusProvedHere]
\label{cor:affine-line-category}
\index{line operators!affine identification}
For $\fg$ simple and $k$ generic, the reduced HT line category (the monodromic braided monoidal category of Corollary~\textup{\ref{cor:braided-category}}), restricted to evaluation modules, is equivalent as a braided monoidal category to $\mathrm{Rep}_q(\fg)$ at $q = e^{i\pi/(k + h^\vee)}$:
\[
 \cC_{\mathrm{line}}^{\mathrm{red}}\big|_{\mathrm{eval}} \;\simeq\; \mathrm{Rep}_q(\fg).
\]
\end{corollary}

\begin{proof}
Theorem~\ref{thm:affine-monodromy-identification}(iii) identifies the braiding and associator. The monoidal structure (tensor product of evaluation modules) is the same on both sides. The equivalence follows from the reconstruction theorem for braided monoidal categories: a braided monoidal category with the same objects, morphisms, associator, and braiding is equivalent.
\end{proof}

\begin{proposition}[RT functor via the negligible-ideal quotient; \ClaimStatusProvedHere]
\label{prop:rt-functor-via-negligible-quotient}
\index{negligible ideal}
\index{Reshetikhin--Turaev functor}
\index{modular tensor category}
\index{Andersen--Paradowski}
Let $\fg$ be a simple Lie algebra, $h^\vee$ its dual Coxeter number, and fix a positive integer level~$k$ with $q = e^{i\pi/(k+h^\vee)}$ a root of unity. Let $\mathrm{Rep}_q^{\mathrm{fus}}(\fg)$ denote the full tensor-ideal generated by tilting modules inside $\mathrm{Rep}_q(\fg)$. Then:
\begin{enumerate}[label=\textup{(\roman*)}]
 \item \emph{(Negligible ideal).} The collection $\cI_{\mathrm{neg}} \subset \mathrm{Rep}_q^{\mathrm{fus}}(\fg)$ of morphisms $f \colon M \to N$ such that $\operatorname{tr}_q(g \circ f) = 0$ for every $g \colon N \to M$ is a tensor ideal. Equivalently, $\cI_{\mathrm{neg}}$ is the ideal generated by negligible objects, those of quantum dimension zero.
 \item \emph{(MTC quotient).} The Verdier-type quotient $\cC^{\mathrm{MTC}}_{k,\fg} := \mathrm{Rep}_q^{\mathrm{fus}}(\fg) / \cI_{\mathrm{neg}}$ is a modular tensor category (MTC), coinciding with the Andersen--Paradowski fusion category at level~$k$ \cite{And92,AP95}; its simple objects are the tilting modules with nonzero quantum dimension, indexed by the fundamental alcove.
 \item \emph{(RT functor).} The Reshetikhin--Turaev construction \cite{RT90} is a ribbon-tangle functor
 \[
  Z^{\mathrm{RT}}_{k,\fg} \colon \mathrm{Rib}(\fg) \longrightarrow \cC^{\mathrm{MTC}}_{k,\fg},
 \]
 equivalently an invariant of framed coloured links valued in $\operatorname{End}(\mathbf{1}) = \CC$. It factors uniquely through the MTC quotient: $Z^{\mathrm{RT}} = \bar{Z} \circ \pi$ with $\pi \colon \mathrm{Rep}_q^{\mathrm{fus}} \twoheadrightarrow \cC^{\mathrm{MTC}}$ and $\bar{Z}$ the induced ribbon functor.
 \item \emph{(Jones specialisation).} For $\fg = \mathfrak{sl}_2$ and the fundamental representation, $Z^{\mathrm{RT}}_{k,\mathfrak{sl}_2}$ on $(1,1)$-tangles (closed framed knots) equals the colored Jones polynomial $J_K(q)$; for $\fg = \mathfrak{sl}_N$ at the defining representation it equals HOMFLYPT at the prescribed specialisation.
 \item \emph{(Bar-complex recovery).} Let $\rho_n^{\mathrm{HT}} \colon B_n \to \operatorname{Aut}(V^{\otimes n})$ be the braid-group representation produced by the bar complex of $V^k(\fg)$ (Theorem~\ref{thm:affine-monodromy-identification}) at generic~$q$. At $q = e^{i\pi/(k+h^\vee)}$ one obtains the corresponding ribbon-category representation of $B_n$ inside $\mathrm{Rep}_q^{\mathrm{fus}}(\fg)$; the RT invariant $Z^{\mathrm{RT}}_{k,\fg}(\hat{\beta})$ of its Markov closure $\hat{\beta}$ equals the composition $\bar{Z} \circ \pi \circ \rho_n^{\mathrm{HT}}(\beta)$ after quantum-trace closure. The quotient $\pi$ is therefore \emph{necessary}: it is what converts the generic-$q$ braided tensor product into the root-of-unity ribbon fusion category underlying Jones.
\end{enumerate}
\end{proposition}

\begin{proof}
(i) That the negligible morphisms form a tensor ideal is a general fact for ribbon categories with a well-defined quantum trace \cite[Prop.~4.4.1]{BaKi01}; closure under tensor product uses $\operatorname{tr}_q(f \otimes g) = \operatorname{tr}_q(f)\operatorname{tr}_q(g)$; closure under composition is immediate from the defining bilinear pairing. The characterisation via quantum-dimension-zero tilting modules is \cite[Thm.~4.1]{And92}.
(ii) Semisimplification of $\mathrm{Rep}_q^{\mathrm{fus}}(\fg)$ by $\cI_{\mathrm{neg}}$ is the Andersen--Paradowski construction \cite{AP95}; modularity (nondegeneracy of the $S$-matrix) is \cite[Thm.~3.3.20]{BaKi01}.
(iii) The RT construction is a ribbon functor out of the ribbon category; it annihilates negligible morphisms because the quantum trace does. Factorisation through $\pi$ is the universal property of Verdier-type quotients by tensor ideals.
(iv) For $\fg = \mathfrak{sl}_2$, fundamental representation, the explicit matching $Z^{\mathrm{RT}} = J_K$ is \cite{RT90}; the $\mathfrak{sl}_N$ case giving HOMFLYPT is \cite[\S3]{Tu94}.
(v) The bar complex produces the braid-group representation on $V^{\otimes n}$ at generic~$q$ by Theorem~\ref{thm:affine-monodromy-identification}(i)--(ii); the analytic continuation $q \to e^{i\pi/(k+h^\vee)}$ lands inside $\mathrm{Rep}_q^{\mathrm{fus}}(\fg)$ (tilting-generated tensor ideal, \cite[Thm.~2.5]{AP95}) because the evaluation modules used are tilting. Quantum-trace closure then factors through $\pi$ by (iii).
\end{proof}

\begin{remark}[The negligible-ideal quotient is not an optional step]
\label{rem:negligible-not-optional}
Without the quotient $\pi$, the generic-$q$ bar-complex braid representation lives in a semisimple braided category that is \emph{not} modular: the non-semisimple extensions at root of unity correspond to tilting modules of quantum dimension zero whose presence destroys nondegeneracy of the $S$-matrix. Computing Jones directly on $\rho_n^{\mathrm{HT}}(\beta)$ before the quotient would produce a well-defined number that is \emph{not} the Jones polynomial in general (it would fail framing anomaly cancellation for non-alcove weights). Corollary~\ref{cor:jones-polynomial} therefore reads: \emph{bar-complex monodromy + negligible-ideal quotient + quantum-trace closure = Jones polynomial}; the middle step is Proposition~\ref{prop:rt-functor-via-negligible-quotient}.
\end{remark}

\begin{corollary}[Jones polynomial from the bar complex; \ClaimStatusProvedHere]
\label{cor:jones-polynomial}
\index{Jones polynomial|textbf}
\index{Kontsevich integral}
For $\fg = \mathfrak{sl}_2$ at positive integer level $k$, the bar complex integrals over $\FM_n(\CC) \times \Conf_n(\RR)$, evaluated on the fundamental weight system of $\mathfrak{sl}_2$, followed by the negligible-ideal quotient $\pi$ of Proposition~\ref{prop:rt-functor-via-negligible-quotient} and quantum-trace closure, compute the colored Jones polynomial $J_K(q)$ of a framed knot $K$ at $q = e^{i\pi/(k+2)}$.
\end{corollary}

\begin{proof}
The argument has three steps: one proved here, one proved in Proposition~\ref{prop:rt-functor-via-negligible-quotient}, one classical.

\medskip
\noindent\textbf{Step 1 (Braid representation from the bar complex).}
The bar complex of $V^k(\mathfrak{sl}_2)$, viewed as a logarithmic $\SCchtop$-algebra, produces a flat connection on configuration spaces by Theorem~\ref{thm:synthesis}. By Theorem~\ref{thm:affine-monodromy-identification}(i)--(ii), on the reduced evaluation comparison surface the monodromy of this connection identifies with the KZ monodromy, and the affine Drinfeld--Kohno theorem compares it with the corresponding braided tensor-product representation of $\mathrm{Rep}_q(\mathfrak{sl}_2)$ at $q = e^{i\pi/(k+2)}$. In particular, the bar complex computes the braid-group representation
\[
 \rho_n^{\mathrm{HT}} = \rho_n^{\mathrm{KZ}} \colon B_n \longrightarrow \Aut(V^{\otimes n})
\]
on tensor powers of any finite-dimensional $\mathfrak{sl}_2$-module~$V$. The integration domain $\FM_n(\CC) \times \Conf_n(\RR)$ parametrises the $E_1$ bar coalgebra: the $\FM_n(\CC)$ factor carries the bar differential and the $\Conf_n(\RR)$ factor carries the $E_1$ coproduct implementing the group-like expansion.

\medskip
\noindent\textbf{Step 2 (Negligible-ideal quotient).}
At $q = e^{i\pi/(k+2)}$ the tilting-generated full tensor-ideal $\mathrm{Rep}_q^{\mathrm{fus}}(\mathfrak{sl}_2)$ carries the negligible ideal $\cI_{\mathrm{neg}}$ of Proposition~\ref{prop:rt-functor-via-negligible-quotient}(i). Passing to the semisimplified quotient $\cC^{\mathrm{MTC}}_{k,\mathfrak{sl}_2} = \mathrm{Rep}_q^{\mathrm{fus}}/\cI_{\mathrm{neg}}$ (Andersen--Paradowski) is necessary for modularity: the $B_n$-action $\rho_n^{\mathrm{HT}}$ descends to the image $\bar\rho_n^{\mathrm{HT}} = \pi \circ \rho_n^{\mathrm{HT}}$ inside the MTC.

\medskip
\noindent\textbf{Step 3 (Braid closure via the RT construction).}
The colored Jones polynomial of a framed knot~$K$, presented as the closure of a braid $\beta \in B_n$, is obtained by applying the ribbon-category quantum trace $\operatorname{tr}_q = \operatorname{tr}(K^{2\rho} \,\cdot\,)$ to $\bar\rho_n^{\mathrm{HT}}(\beta)$ in $\cC^{\mathrm{MTC}}_{k,\mathfrak{sl}_2}$:
\[
 J_K(q) \;=\; \operatorname{tr}_q\bigl(\bar\rho_n^{\mathrm{HT}}(\beta)\bigr)
 \;=\; \operatorname{tr}_q\bigl(\pi \circ \rho_n^{\mathrm{HT}}(\beta)\bigr).
\]
This is the Reshetikhin--Turaev construction \cite{RT90} applied to the MTC of Proposition~\ref{prop:rt-functor-via-negligible-quotient}(ii)--(iii): $\cC^{\mathrm{MTC}}_{k,\mathfrak{sl}_2}$ is modular, its ribbon element implements framing, the quantum trace implements Markov closure.

\medskip
\noindent\textbf{Conclusion.}
Composing Steps~1--3: bar-complex monodromy, negligible-ideal quotient, quantum-trace closure, all three steps constructive and each proved (Steps~1 and~2 here and in Proposition~\ref{prop:rt-functor-via-negligible-quotient}; Step~3 classical). The output is $J_K(q)$.
\end{proof}

\begin{remark}[Scope: $\mathfrak{sl}_N$, exceptional types, and the universal celestial analogue]
\label{rem:rt-scope}
Proposition~\ref{prop:rt-functor-via-negligible-quotient} and Corollary~\ref{cor:jones-polynomial} extend verbatim to $\fg = \mathfrak{sl}_N$ (HOMFLYPT at the defining representation) and to the classical series $\mathfrak{so}_N$, $\mathfrak{sp}_{2N}$ (two-variable Kauffman) using the Andersen--Paradowski MTC at root of unity \cite{AP95}. For exceptional types $G_2, F_4, E_{6,7,8}$, the MTC is also constructed by Andersen--Paradowski and gives the corresponding quantum invariants. An analogous negligible-ideal quotient governs the universal celestial holography programme (Theorem~\ref{thm:affine-monodromy-identification}, extended to $\mathcal{W}$-algebras via Remark~\ref{rem:affine-scope}(b)); the twistor-gravity quotient is the $\mathcal{W}$-algebra analogue of $\cI_{\mathrm{neg}}$, semisimplifying the $\mathcal{W}$-Yangian's generic-$q$ braided tensor product. Making this explicit is Conjecture~\ref{conj:rmatrix}.
\end{remark}

\begin{corollary}[First proved holographic dictionary; \ClaimStatusProvedHere]
\label{cor:holographic-dictionary}
\index{holographic dictionary!affine}
The affine lineage provides the first unconditionally proved instance of the full holographic dictionary of the Swiss-cheese programme:
\begin{center}
\renewcommand{\arraystretch}{1.3}
\begin{tabular}{lll}
\textbf{3d HT bulk} & \textbf{2d boundary} & \textbf{Line operators} \\
\hline
Affine PV sigma model & $V^{k+\hbar\delta_\fg}(\fg)$ & $\mathrm{Rep}_q(\fg)$ \\
Bar differential on $\FM_k(\CC)$ & Chiral OPE & Braiding \\
Deconcatenation on $\Conf_k^<(\RR)$ & Factorisation coproduct & Coalgebra structure \\
Reduced HT associator & KZ associator & Drinfeld associator \\
\end{tabular}
\end{center}
The bulk-to-boundary functor is the half-space quantization of Theorem~\textup{\ref{thm:affine-half-space-bv}}. The boundary-to-line functor is open-colour Koszul duality: $\cA^!_{\mathrm{line}} \simeq Y_\hbar(\fg)$. The identification of the line category with the quantum group representation category is Theorem~\textup{\ref{thm:affine-monodromy-identification}}.
\end{corollary}

\begin{remark}[Scope and frontier]
\label{rem:affine-scope}
Theorem~\ref{thm:affine-monodromy-identification} proves the monodromy identification for evaluation modules of the affine lineage. Three natural extensions remain:
\begin{enumerate}[label=\textup{(\alph*)}]
 \item \emph{Beyond evaluation modules}: for infinite-dimensional modules in category $\cO_k(\widehat{\fg})$, the identification $\rho^{\mathrm{HT}} = \rho^{\mathrm{KZ}}$ requires additional convergence analysis of the $\h$-adic retract (Proposition~\ref{prop:analytic-retract}). The DK-2 equivalence $\cO_k^{\mathrm{int}}(\widehat{\fg}) \simeq \mathrm{Rep}_q(\fg)$ of Kazhdan--Lusztig~\cite{KL93} suggests this should hold.
 \item \emph{Beyond the affine lineage}: for $\mathcal{W}$-algebras obtained by Drinfeld--Sokolov reduction, the Swiss-cheese transferred operations $m_k^{\mathrm{SC}}$ ($k \geq 3$) do not vanish. The one-loop collapse mechanism breaks, and the full $A_\infty$ superconnection contributes. The reduced monodromy is then the $\mathcal{W}$-algebra $R$-matrix, which should equal the spectral $R$-matrix of the $\mathcal{W}$-Yangian; this is the content of Conjecture~\ref{conj:rmatrix} in the non-formal Swiss-cheese regime.
 \item \emph{Higher genus}: at genus $g \geq 1$, the KZ connection is replaced by the KZB connection (Computation~\ref{comp:genus1-kzb}), and the monodromy acquires modular corrections controlled by the curvature $\kappa(\cA) \cdot \omega_g$. The identification should lift to elliptic and higher-genus $R$-matrices, connecting to the chromatic tower (Conjecture~\ref{conj:chromatic}).
\end{enumerate}
\end{remark}

% Bibliography entries consolidated into main.tex.

\section{Synthesis}
\label{sec:loght-core-supplement}

\begin{remark}[Logarithmic HT monodromy and the Humbert divisor]
\label{rem:loght-core-humbert}
The logarithmic HT-monodromy framework specialises at Mukai-enhanced K3 to the holonomic $\mathcal{D}$-module $\mathcal{L}^{\Delta_5}$ on $\overline{\mathcal{A}_2}$ with regular singularities on the Humbert divisors $H_1 \cup H_4$. The monodromy of $\mathcal{L}^{\Delta_5}$ around $H_1$ has order $8$, identified via Bruinier $2002$ Proposition~$5.1$ Heegner Chern-class reciprocity with the Mukai-doubling Koszul conductor $K^{\kappa_{\mathrm{ch}}} = 2\,c_+(\mathrm{Mukai}(K3)) = 8$ and the Lusztig specialisation $\ell = 8$ at $\zeta^8 = 1$, $\hbar^2 = -1/8$. Universal identity $\hbar^2 \cdot K^{\kappa_{\mathrm{ch}}} = -1$ on the Lorentzian-lattice-parametric $\mathcal{B}$-family.
\end{remark}
